

% ----------------------------------------------------------------------
%                   LATEX TEMPLATE FOR PhD THESIS
% ----------------------------------------------------------------------

% based on Harish Bhanderi's PhD/MPhil template, then Uni Cambridge
% http://www-h.eng.cam.ac.uk/help/tpl/textprocessing/ThesisStyle/
% corrected and extended in 2007 by Jakob Suckale, then MPI-CBG PhD programme
% and made available through OpenWetWare.org - the free biology wiki


%: Style file for Latex
% Most style definitions are in the external file PhDthesisPSnPDF.
% In this template package, it can be found in ./Latex/Classes/
\documentclass[twoside, 12pt]{thesis}


%: Macro file for Latex
% Macros help you summarise frequently repeated Latex commands.
% Here, they are placed in an external file /Latex/Macros/MacroFile1.tex
% An macro that you may use frequently is the figuremacro (see introduction.tex)
\newcommand{\comment}[1]{}
\newcommand{\myparagraph}[1]{\paragraph{#1}\mbox{}\\}
\newcommand{\pw}[4]{
	\begin{cases}\displaystyle
		#1 & \mbox{#2},\\ 
		#3 & \mbox{#4}
	\end{cases}
}

\newcommand{\cref}[1]{Chap.~\ref{#1}}
\newcommand{\Cref}[1]{Chapter~\ref{#1}}
\newcommand{\sref}[1]{Sec.~\ref{#1}}
\newcommand{\Sref}[1]{Section~\ref{#1}}
\newcommand{\aref}[1]{App.~\ref{#1}}
\newcommand{\Aref}[1]{Appendix~\ref{#1}}
\newcommand{\fref}[1]{Fig.~\ref{#1}}
\newcommand{\Fref}[1]{Figure~\ref{#1}}
\newcommand{\tref}[1]{Tab.~\ref{#1}}
\newcommand{\Tref}[1]{Table~\ref{#1}}
\newcommand{\pgref}[1]{p.~\pageref{#1}}
\newcommand{\Pgref}[1]{page~\pageref{#1}}



\def\eqrefs#1#2{(\protect\ref{#1}--\protect\ref{#2})}

\newcommand{\alert}[1]{\textcolor{red}{#1}}
\newcommand{\PN}[1][N]{\ensuremath{\text{P}_{#1}}}
\newcommand{\MCPN}[1][N]{\ensuremath{\text{MCP}_{#1}}}
\newcommand{\SN}[1][N]{\ensuremath{\text{S}_{#1}}}
\newcommand{\SPN}[1][N]{\ensuremath{\text{SP}_{#1}}}
\newcommand{\Oh}{\ensuremath{\mathcal{O}}}

\newcommand{\R}[1][]{\mathbb{R}^{#1}}
\newcommand{\mat}[1]{\mathbf{#1}}
\newcommand{\src}{q}

\newcommand{\sla}{\ensuremath{\shortleftarrow}}
\newcommand{\sra}{\ensuremath{\shortrightarrow}}

\newcommand{\br}{\ensuremath{\mathbf{r}}}
\newcommand{\bv}{\ensuremath{\mathbf{v}}}
\newcommand{\bn}{\ensuremath{\mathbf{n}}}
\newcommand{\be}{\ensuremath{\mathbf{e}}}
\newcommand{\bJ}{\ensuremath{\mathbf{J}}}

\newcommand{\bomega}{\mathbf{\Omega}}
\newcommand{\Sphere}{\ensuremath{S_2}}
\newcommand{\Emin}{\ensuremath{E_{\text{min}}}}
\newcommand{\Emax}{\ensuremath{E_{\text{max}}}}

\newcommand{\der}[3][]{\ensuremath{\frac{\text{d}^{#1}{#2}}{\text{d} {#3}^{#1}}}}
\newcommand{\pd}[3][]{\ensuremath{\frac{\partial^{#1}#2}{\partial #3^{#1}}}}
\newcommand{\grad}{\nabla}
\newcommand{\gradcomp}[1]{\nabla_{\!#1}}
\newcommand{\lapv}{\nabla^2}
\newcommand{\lap}{\nabla^2}

\renewcommand{\d}[1]{\ensuremath{\text{d}#1\,}}
\newcommand{\intE}[4][]{\int_{#2}^{#3}#4\,\d{E#1}}
\newcommand{\intEg}[3]{\int_{E^{#1}}^{E^{#2}}#3\,\d{E}}
\newcommand{\intA}[2][]{\int_{\Sphere}#2\,\d{\bomega#1}}
\newcommand{\intX}[2][X]{\int_{#1}#2\,\d{x}}

\newcommand{\sumg}{\sum_{{g'\neq g}}}
\newcommand{\summ}[2]{\sum_{#1}^{#2}}
\newcommand{\suma}[3][l]{\summ{#1=#2}{#3}}

\newcommand{\abs}[1]{\ensuremath{\left\vert#1\right\vert}}
\newcommand{\norm}[2][]{\ensuremath{\lVert#2\rVert_{#1}}}

\newcommand{\pX}[1][]{\partial X^{#1}}
\newcommand{\pV}[1][]{\partial V}
\newcommand{\VV}{V}

\newcommand{\polar}{\vartheta}
\newcommand{\azimuthal}{\varphi}
\newcommand{\sint}{\sin \polar}
\newcommand{\sinp}{\sin \azimuthal}
\newcommand{\cost}{\cos \polar}
\newcommand{\cosp}{\cos \azimuthal}
\newcommand{\db}{\ensuremath{\d\xi}}

\newcommand{\Lp}[1][p]{L^{#1}}
\newcommand{\Hp}[1][p]{H^{#1}}
\newcommand{\Id}{\ensuremath{I}}
\newcommand{\Dom}[1]{\mathcal{D}(#1)}

\newcommand{\Y}[2]{Y_{#1}^{#2}}
\renewcommand{\P}[2]{P_{#1}^{#2}}
\newcommand{\YY}[1][n]{\mathbf{Y}_{#1}}
\newcommand{\Yc}[2]{\bar{Y}_{#1}^{#2}}
\newcommand{\YYc}[1][n]{\bar{\mathbf{Y}}_{#1}}

\newcommand{\PP}[1][n]{\mathbb{P}^{(#1)}}
\newcommand{\PPP}[2][n]{P^{(#1)}_{#2}}
\renewcommand{\AA}[1][n]{\mathbb{A}^{(#1)}}
\newcommand{\AAA}[2][n]{A^{(#1)}_{#2}}
\newcommand{\BB}[1][n]{\mathbb{B}^{(#1)}}
\newcommand{\BBB}[2][n]{B^{(#1)}_{#2}}
\newcommand{\CC}[1][n]{\mathbb{C}^{(#1)}}
\newcommand{\CCC}[2][n]{C^{(#1)}_{#2}}


\newcommand{\kron}[2]{\delta_{#1#2}}
\newcommand{\Sym}[1]{\left[#1\right]_{\mbox{\scriptsize sym}}}
\newcommand{\eye}{\mathbb{I}}
\newcommand{\detr}[1][n]{\mathscr{D}_{#1}}
\newcommand{\contr}[1][n]{\cdot}%\underset{#1}{\cdot}}
\newcommand{\perm}[1]{\pi(#1)}

\newcommand{\angmom}[2]{\ensuremath{\phi_{#1}^{#2}}}
\newcommand{\fmom}[1][n]{\Phi_{#1}}
\newcommand{\qmom}[1][n]{Q_{#1}}
\newcommand{\Sa}[2][]{\Sigma^{#1}_{#2}}
\newcommand{\mst}{\underline{\sigma_t}}
%\newcommand{\Sym}[2]{\mathscr{S}^{(#1)}\left[#2\right]}

\newcommand{\angflux}{\ensuremath{\psi}}

\DeclareMathOperator{\sgn}{sgn}
\DeclareMathOperator{\esssup}{ess\,sup\,}
\DeclareMathOperator{\argmin}{argmin\,}
\DeclareMathOperator{\diag}{diag\,}
\DeclareMathOperator{\col}{col\,}
\DeclareMathOperator{\trace}{tr\,}

\usepackage{colortbl}
\newlength\celldim \newlength\fontheight \newlength\extraheight
\setlength\celldim{2em}
\settoheight\fontheight{A}
\setlength\extraheight{\celldim - \fontheight}

\newcolumntype{A}{ >{\centering } p{1em}  <{\rule[-.5\extraheight]{0pt}{\extraheight+\fontheight}}}
\newcolumntype{B}{ >{\centering } p{11em}  <{\rule[-.5\extraheight]{0pt}{\extraheight+\fontheight}}}
\newcolumntype{C}{ >{\centering } p{11em} <{\rule[-.5\extraheight]{0pt}{\extraheight+\fontheight}}}
\newcolumntype{D}{ >{\centering } p{4em}  <{\rule[-.5\extraheight]{0pt}{\extraheight+\fontheight}}}
\newcommand{\nl}{\tabularnewline\hline}




%: ----------------------------------------------------------------------
%:                  TITLE PAGE: name, degree,..
% ----------------------------------------------------------------------
% below is to generate the title page with crest and author name

\title{Mathematical Modeling\\ of Neutron Transport}



% ----------------------------------------------------------------------
% The section below defines www links/email for author and institutions
% They will appear on the title page of the PDF and can be clicked
\ifpdf
  \author{\href{mailto:mhanus@kma.zcu.cz}{Milan Hanu{\v s}}}
%  \cityofbirth{born in XYZ} % uncomment this if your university requires this
%  % If city of birth is required, also uncomment 2 sections in PhDthesisPSnPDF
%  % Just search for the "city" and you'll find them.
  \collegeordept{\href{http://www.kma.zcu.cz}{Department of Mathematics}}
  \university{\href{http://www.zcu.cz}{University of West Bohemia, Pilsen}}

  % The crest is a graphics file of the logo of your research institution.
  % Place it in ./0_frontmatter/figures and specify the width
  \crest{\includegraphics[width=4cm]{logo}}
  
% If you are not creating a PDF then use the following. The default is PDF.
\else
  \author{Milan Hanu{\v s}}
  \collegeordept{Department of Mathematics}
  \university{University of West Bohemia, Pilsen}
  \crest{\includegraphics[width=4cm]{logo}}
\fi

\renewcommand{\submittedtext}{Thesis}
\supervisor{Doc. Ing. Marek Brandner, Ph.D.}
\degree{Doctor of Philosophy (Ph.D.)}
\degreedate{07/2012}


% ----------------------------------------------------------------------
       
% turn of those nasty overfull and underfull hboxes
\hbadness=10000
\hfuzz=50pt

 
%: --------------------------------------------------------------
%:                  FRONT MATTER: dedications, abstract,..
% --------------------------------------------------------------

\begin{document}

%\language{english}

% sets line spacing
\renewcommand\baselinestretch{1.2}
\baselineskip=18pt plus1pt


%: ----------------------- generate cover page ------------------------

 

% command to print the title page with above variables
\maketitle

%: ----------------------- cover page back side ------------------------
% Your research institution may require reviewer names, etc.
% This cover back side is required by Dresden Med Fac; uncomment if needed.

%\newpage
%\vspace{10mm}
%1. Reviewer: Name

%\vspace{10mm}
%2. Reviewer: 

%\vspace{20mm}
%Day of the defense:

%\vspace{20mm}
%\hspace{70mm}Signature from head of PhD committee:



%: ----------------------- abstract ------------------------

% Your institution may have specific regulations if you need an abstract and where it is to be placed in the document. The default here is just after title.


% Thesis Abstract -----------------------------------------------------


%\begin{abstractslong}    %uncommenting this line, gives a different abstract heading
\begin{abstracts}        %this creates the heading for the abstract page

Put your abstract or summary here, if your university requires it.

\end{abstracts}
%\end{abstractlongs}


% ---------------------------------------------------------------------- 


% The original template provides and abstractseparate environment, if your institution requires them to be separate. I think it's easier to print the abstract from the complete thesis by restricting printing to the relevant page.
% \begin{abstractseparate}
%   
% Thesis Abstract -----------------------------------------------------

\begin{alwayssingle} \pagestyle{empty}
  \begin{center}
  \vspace*{1.5cm}
  {\Large \bfseries  Abstract}
  \end{center}
  \vspace{0.5cm}
\begin{quote}

The subject of this work is computational modeling of neutron transport relevant to economical and safe operation of
nuclear facilities. The general mathematical model of neutron transport is provided by the linear Boltzmann's transport
equation\comment{This first-order partial differential equation with integral terms describes the dependence of neutron
field on six phase-space variables (position, direction, energy)} and the thesis begins with its precise mathematical
formulation and presentation of known conditions for its well-posedness.

\comment{High-dimensionality and complicated structure of the equation for real-world problems preclude analytical
solution and present serious difficulties when a computational solution is attempted. Several dimension-reduction
techniques thus need to be employed along with appropriate numerical solution schemes. This is the subject of the}
In the following part, we study approximation methods for the transport equation, starting with the classical
discretization of energetic dependence and followed by the review of two most widely used methods for approximating 
directional dependence (the $\SN$ and $\PN$ methods). While these methods are usually presented independently of each 
other, we show that they can be put into a single framework of Hilbert space projection techniques. This fact is then
used in conjunction with the results of the first part to rigorously prove rotational invariance of the $\PN$
equations and to analyze convergence of the basic iterative scheme for solving the $\SN$ equations. This part of the
thesis is concluded by the description of a finite element method for the final discretization of spatial dependence and a discussion of the solution of the resulting system of
algebraic equations.

The main new results are contained in the following two chapters focusing on the simplified $\PN$ approximation, which
is a computationally more convenient albeit not as mathematically well-founded variant of the $\PN$ approximation. We
prove well-posedness of the weak form of the $\SPN[3-7]$ equations and present a new way of deriving the equations from
an alternative set to the $\PN$ equations, obtained from special linear combination of
spherical harmonics.% (the so-called Maxwell-Cartesian spherical harmonics, hence the abbreviation $\MCPN$ for this
% approximation), u%tilizing its tensor structure. Algebraic manipulations allowing to obtain the $\SPN[3]$-equivalent equations from the
%$\MCPN[3]$ set in an interior of a homogeneous region are explicitly given in \sref{sec:mcp3_red}.

The final part of the thesis contains numerical examples of the $\SN$ and $hp$-adaptive $\SPN$ calculations using a
neutronics framework that has been implemented by the author to the $hp$-adaptive finite element library Hermes2D. The $\SPN[1]$
(or diffusion) model also serves as a basis of a real-world reactor calculation suite co-developed by the author for the
purposes of ``Project TA01020352 -- Increasing utilization of nuclear fuel through optimization of an inner fuel cycle and
calculation of neutron-physics characteristics of nuclear reactor cores''. An example benchmark used to test the code
concludes the thesis.
 

\vspace*{1cm}
%\begin{quote}
{\large \bfseries  Keywords:}

The neutron diffusion approximation, whereby the NTE is reduced to a second-order elliptic PDE (or, when energy
dependence is taken into account implicitly, a weakly coupled non-symmetric system of second-order PDEs with
positive-definite symmetric part -- the so-called \textit{multigroup neutron diffusion approximation}), also forms the

\end{quote}

\clearpage
\selectlanguage{czech}
  \begin{center}
  \vspace*{1.5cm}
  {\Large \bfseries  Abstrakt}
  \end{center}
  \vspace{0.5cm}
\begin{quote}

Práce se zabývá matematickým a numerickým modelováním transportu neutronů, se zaměřením na výpočty neutronových
charakteristik jaderných reaktorů. Obecný matematický model transportu neutronů je reprezentován lineární Boltzmannovou
transportní rovnicí. Práce začíná její přesnou matematickou formulaci a přehledem\linebreak výsledků týkajících se její
řešitelnosti ve druhé kapitole. Následující kapitoly jsou zaměřeny na přibližné metody řešení této rovnice.

Po stručném popisu klasické diskretizace energetické závislosti je hlavní část třetí kapitoly věnována aproximaci
směrové závislosti pomocí dvou stěžejních metod -- metody diskrétních ordinát ($\SN$) a metody sférických harmonických
funkcí ($\PN$). Zatímco obvykle jsou tyto metody formulovány nezávisle, v práci je ukázáno, jak je lze obě popsat pomocí
jednotného rámce jako projekci na podprostor Hilbertova prostoru funkcí definovaných na sféře. Této skutečnosti je
posléze využito při důkazu rotační invariantnosti $\PN$ rovnic a při konvergenční analýze základní iterační
metody pro řešení $\SN$ soustavy.
Třetí kapitola je zakončena popisem aplikace metody konečných prvků na finální diskretizaci prostorové závislosti.

Hlavní nové výsledky této práce se týkají metody zjednodušených sférických harmonických funkcí ($\SPN$), jež představuje
výpočetně efektivní aproximaci metody $\PN$. Ve čtvrté kapitole je standardním \linebreak způsobem odvozena slabá
formulace $\SPN$ rovnic a dokázána její korektnost pro $N = 3,5,7$. V páté kapitole je pak odvozena
nová soustava parciálních diferenciálních rovnic odpovídající $\PN$ aproximaci ($\MCPN$ aproximace). Na příkladu $\MCPN[3]$ aproximace je ukázáno, jak lze využít
tenzorovou strukturu těchto rovnic k transformaci na soustavu ekvivalentní s $\SPN[3]$ aproximací.

V šesté kapitole je popsána implementace $\SN$ a $\SPN$ aproximací do knihovny Hermes2D a na několika
příkladech ukázány základní vlastnosti těchto aproximací.  Speciální pozornost je věnována implementaci nespojité
Galerkinovy metody (pro $\SN$ aproximaci)%, umožňující využít speciální způsob asemblace soustav PDR v knihovně
% Hermes2D, v níž je každá neznámá funkce aproximována na vlastním konečně-prvkovém prostoru. 
a modifikaci standardního indikátoru chyby pro
$hp$-adaptivitu v Hermes2D pro $\SPN$ aproximaci. Práce je ukončena ukázkou řešení standardního 3D benchmarku pomocí
mnohagrupového difúzního kódu, který autor na základě zkušeností s vývojem neutronických modulů v knihovně Hermes2D
vyvinul pro účely projektu ``TA01020352 -- Zvýšení využití jaderného paliva pomocí optimalizace vnitřního palivového
cyklu a výpočtu neutronově-fyzikálních charakt. aktivních zón jaderných reaktorů''.



  \vspace*{1cm}
{\large \bfseries  Klíčová slova:}

%\begin{quote}

The neutron diffusion approximation, whereby the NTE is reduced to a second-order elliptic PDE (or, when energy
dependence is taken into account implicitly, a weakly coupled non-symmetric system of second-order PDEs with
positive-definite symmetric part -- the so-called \textit{multigroup neutron diffusion approximation}), also forms the


\end{quote}

\end{alwayssingle}\selectlanguage{english}


% ---------------------------------------------------------------------- 

% \end{abstractseparate}


%: ----------------------- tie in front matter ------------------------

\frontmatter
% Thesis Dedictation ---------------------------------------------------

\begin{dedication} %this creates the heading for the dedication page

To ...

\end{dedication}

% ----------------------------------------------------------------------
% Thesis Acknowledgements ------------------------------------------------


%\begin{acknowledgementslong} %uncommenting this line, gives a different acknowledgements heading
\begin{acknowledgements}      %this creates the heading for the acknowlegments

I would like to acknowledge the thousands of individuals who have coded for the LaTeX project for free. It is due to their efforts that we can generate professionally typeset PDFs now.

\end{acknowledgements}
%\end{acknowledgmentslong}

% ------------------------------------------------------------------------





%: ----------------------- contents ------------------------

\setcounter{secnumdepth}{3} % organisational level that receives a numbers
\setcounter{tocdepth}{2}    % print table of contents for level 3
\tableofcontents            % print the table of contents
% levels are: 0 - chapter, 1 - section, 2 - subsection, 3 - subsection


%: ----------------------- list of figures/tables/todos ------------------------

\listoffigures	% print list of figures
\listoftables  % print list of tables

%\todototoc
%\listoftodos


%: ----------------------- glossary ------------------------

% Tie in external source file for definitions: /0_frontmatter/glossary.tex
% Glossary entries can also be defined in the main text. See glossary.tex
% this file is called up by thesis.tex
% content in this file will be fed into the main document

% Glossary entries are defined with the command \nomenclature{1}{2}
% 1 = Entry name, e.g. abbreviation; 2 = Explanation
% You can place all explanations in this separate file or declare them in the middle of the text. Either way they will be collected in the glossary.

% required to print nomenclature name to page header
\markboth{\MakeUppercase{\nomname}}{\MakeUppercase{\nomname}}


% ----------------------- contents from here ------------------------

% chemicals
\nomenclature{DAPI}{4',6-diamidino-2-phenylindole; a fluorescent stain that binds strongly to DNA and serves to marks the nucleus in fluorescence microscopy} 
\nomenclature{DEPC}{diethyl-pyro-carbonate; used to remove RNA-degrading enzymes (RNAases) from water and laboratory utensils}
\nomenclature{DMSO}{dimethyl sulfoxide; organic solvent, readily passes through skin, cryoprotectant in cell culture}
\nomenclature{EDTA}{Ethylene-diamine-tetraacetic acid; a chelating (two-pronged) molecule used to sequester most divalent (or trivalent) metal ions, such as calcium (Ca$^{2+}$) and magnesium (Mg$^{2+}$), copper (Cu$^{2+}$), or iron (Fe$^{2+}$ / Fe$^{3+}$)}



 

%\begin{multicols}{2} % \begin{multicols}{#columns}[header text][space]
%\begin{footnotesize} % scriptsize(7) < footnotesize(8) < small (9) < normal (10)

\printnomenclature[1.5cm] % [] = distance between entry and description
\label{nom} % target name for links to glossary
    
%\end{footnotesize}
%\end{multicols}



%: --------------------------------------------------------------
%:                  MAIN DOCUMENT SECTION
% --------------------------------------------------------------

% the main text starts here with the introduction, 1st chapter,...
\mainmatter

\renewcommand{\chaptername}{} % uncomment to print only "1" not "Chapter 1"


%: ----------------------- subdocuments ------------------------

  \chapter{Introduction}

Computer simulation of radiative transfer of energy is an important task in many engineering and research areas, as
diverse as biomedicine, astrophysics, optics or nuclear engineering. In nuclear engineering, the area of primary
interest in this work, there are two main goals of evaluating computer models of radiative transfer. The first is to
simulate short-term transient behavior of nuclear devices under given initial conditions such as geometry and material
configuration. The second is to determine under which conditions such devices (in this case typically nuclear reactor
cores) will be capable of long-term, stable operation satisfying certain safety, technical and economical limitations,
with only a minimal human intervention. Repeated calculations of the second type form the basis for designing new
nuclear reactors or optimizing fuel reloading of existing ones. Optimization of fuel reloading schemes for nuclear
reactors is the topic of a major research and development project that is being investigated at author's department
\footnote{Project TA01020352 -- Increasing utilization of nuclear fuel through optimization of an inner fuel cycle and
calculation of neutron-physics characteristics of nuclear reactor cores. Principal investigators: R. {\v C}ada
(University of West Bohemia) and J. Rataj (Czech Technical University). \label{ftn:TACR}}.
Author's participation in this project during the course of his doctoral studies involved the development of a
neutron-physical calculation module, employed by the overall optimization suite to evaluate fitness of its candidate
configurations. This fact largely influenced the direction for development and analysis of mathematical models and
numerical methods described in this thesis.

All forms of radiative transfer are governed by physical laws that can be abstracted by a single mathematical model,
represented by the Boltzmann transport equation. In nuclear reactors, the effects of neutron-induced reactions dominate
those caused by other types of particles and we will therefore consider the transport equation for neutrons (shortly
NTE\nomenclature[X]{neutron transport equation}) in this thesis, even though it has the same form for other types of
non-charged particles, such as photons. While the short-term transient simulations require accurate solution methods for
the time-dependent NTE, a quasi-steady state solution (a sequence of steady state calculations) is generally sufficient
to capture slow changes in core configuration and material characteristics during its long-term stable operation.
As we focus on the latter application domain in this thesis, its main subject of study will be the \textit{steady
state neutron transport equation}, shortly NTE. In the final chapter, we will consider the quasi-steady state NTE
(i.e., without time derivative but with time-dependent coefficients) together with the most important feedbacks from
varying thermal/hydraulic conditions in the core and fuel depletion.

We will introduce the NTE in Chap. \ref{chap:nte-review} as an integro-differential equation in 6 dimensions. This high
dimensionality requires either the use of direct simulation of neutrons and using statistical methods to obtain the
required physical quantities (the Monte Carlo approach) or a deterministic approach involving multiple discretizations.
As the second approach is still preferable in terms of overall efficiency, we choose it as a basis for our methods and
review the classical discretization methods in the second part of Chap. \ref{chap:nte-review}. The transport equation will
then be presented in more detail in Chap. \ref{chap:NTE}, together with the review of existing solvability results. 

The following subsections present the main results of author's work during his Ph.D. studies.

\paragraph{The $\MCPN$ approximation and its relation to the $\SPN$ approximation} 
One of the traditional discretization methods that converts the 6-dimensional integro-differential NTE into a system of
familiar 3-dimensional partial differential equations (PDEs\nomenclature[Z]{PDE}{Partial differential equation}) is the
\textit{method of spherical harmonics}, denoted $\PN$\nomenclature[Z]{\PN}{Method of spherical harmonics of order $N$} 
\nomenclature[Z]{\SN}{Method of discrete ordinates of order $N$}. This method is recalled in Chap.
\ref{chap:nte-review} and a new alternative set (equivalent to the original $\PN$ set) is developed in Chap.
\ref{chap:MCSH}. From this set, it is possible to directly derive another traditional approximation of neutron transport
-- the \textit{simplified spherical harmonic method}, or $\SPN$ \nomenclature[Z]{\SPN}{Method of simplified spherical
harmonics of order $N$} -- which is the subject of the second part of Chap. \ref{chap:MCSH}. So far, the author has only
been able to derive in this way the $\SPN$ equations in an interior of a homogeneous region, so in order to obtain a
practical solution method, the usual $\SPN$ interface and boundary conditions will be applied. As a result, the standard
$\SPN$ method is obtained, so the contribution of Chap. \ref{chap:MCSH} is rather theoretical, showing a new way of
arriving at the $\SPN$ equations (and an alternative form of the classical $\PN$ equations).

The $\SPN$ method (particularly the $\SPN[3]$) already simplifies the NTE to the extent that it is applicable to
day-by-day whole-core calculations on usual workstations with a few computational cores or small-scale parallel machines
with tens to a few hundred cores, which are typical machines available to nuclear engineering
companies\footnote{\label{sjsexp}from personal experience of the author coming out of the long-term collaboration with
the Czech nuclear engineering company {\v S}koda JS; more generally, see the discussion in \cite[Sec.
2.4]{Sanchez7}}. As is well known and will be recalled in \ref{chap:SPN}, when the $\SPN[3]$ method is applied to the
typical reactor core calculations, its solution captures most of the features of the true solution of the NTE. Combined
with its efficiency that allows this method to be used ``of the desk'' (without the need of submitting the job to some
supercomputing center, waiting for it to come to the front of an execution queue and gathering the results) makes it
attractive for physicists to quickly test their empirical approximations used throughout their production code, which is
usually based on the strongest transport approximation -- the diffusion approximation.

\myparagraph{Finite element framework for 2D neutron diffusion or transport problems}
The neutron diffusion approximation, whereby the NTE is reduced to a second-order elliptic PDE (or, when energy
dependence is taken into account implicitly, a weakly coupled non-symmetric system of second-order PDEs with
positive-definite symmetric part -- the so-called \textit{multigroup neutron diffusion approximation}), also forms the
basis of our neutronics module for the optimization code\footnote{Note that the $\SPN[1]$ model is actually equivalent
to the neutron diffusion equation.}
However, to obtain the final discrete algebraic system of equations, it uses the finite element method. 
This distinguishes it from the majority of other codes used for similar purposes, which are usually based on the so-called \textit{nodal method}\footnote{See e.g.
\cite{opt1,opt2,opt3} for the specific application area of core reloading optimization; some other nodal codes widely
employed in various whole-core calculations are tabulated in \cite{mox-bench}. One of the two Czech standardized codes
for reactor calculations -- ANDREA \cite{ANDREA} -- also falls into this category (the other -- MOBY-DICK
\cite{MOBYDICK} -- is based on a finite-difference method on a structured mesh)\label{ftn:nodal}}.
Generally speaking, a nodal method is a coarse mesh finite volume method iteratively combined with fine-level correction steps, specially tailored to the
neutron diffusion (or recently $\SPN$) model and reactor core domain (i.e., typically, coarse level cells
correspond to real fuel assemblies and the correction consists of analytic solution of the diffusion equation in a
geometrically simple homogeneous region). Its advantage is the speed and overall efficiency, but it is greatly limited
in geometrical flexibility. Because of the way the nodal equations are derived, it also requires the
\textit{homogenization} procedure to represent each coarse cell by a single set of material coefficients and the 
corresponding \textit{dehomogenization} procedure to reconstruct the fine structure of the solution needed for further 
computations\footnote{We note that in our verification examples in Chap. \alert{section}, we also use homogenized input 
data and we expect such data to be used also during the optimization runs for the sake of efficiency; however, it is not
the \textit{requirement} of the code.} (where the latter, in particular, is difficult to formulate in general cases).
Moreover, the convergence and stability of the method is, to the knowledge of the author, not very well understood 
(\cite{ZiminComm}).

Although the author also participates in the development of a nodal module for the code MOBY-DICK (see footnote
\ref{ftn:nodal}; \cite{Hanus2}, \cite[Chap. 4]{Hanus3}), only the finite-element approximation of both the diffusion and
the $\SPN$ model will be considered in this thesis. In fact, one of the main objectives of the thesis is to show the
feasibility of using finite element approximation in a quasi-steady state whole-core calculation code. In order to
provide a viable alternative to coarse nodal and structured finite-difference discretizations, which often incorporate
hard-coded formulas also for the solutions of the associated discrete algebraic systems, a need for automatic mesh
adaptivity and fast algebraic solvers has been recognized from the beginning. Rather than developing an adaptive finite
element code from scratch, it has been decided to modify an existing library for the needs of neutron transport
modelling. This led the author to the finite-element C++ framework Hermes (\cite{Hermes-project}) with advanced
\textit{hp-adaptivity} capabilities (\cite{Hermes-hanging-nodes}) and a unique way of assembling coupled systems of PDEs (\cite{Hermes-thermoelasticity}).
Author's main contributions to the Hermes project involved:
\begin{itemize}
    \item development of an interface for various existing sparse,
direct and iterative algebraic solvers (which also required reworking the CMake build system of Hermes),
	\item development of a multigroup neutron diffusion framework, simplifying and unifying the formulation of
	multiregion, multigroup neutron diffusion problems within Hermes,
	\item development of the discontinuous Galerkin framework (together with L. Korous, the main developer of Hermes at
	present time) and
	\item extension of the h-adaptivity capabilities by the standard a-posteriori error estimation for elliptic problems (which
involves solution jumps over element interfaces and thus uses elements of the discontinuous Galerkin framework). 
\end{itemize}
More details on these contributions will be given in Chap. \alert{ref}, where also the extension of the multigroup
neutron diffusion framework to include the $\SPN$ model (for orders $N \leq 9$, but easily extendable to an arbitrary
order) will be presented, together with specially tailored error estimate to drive the hp-adaptivity of Hermes. To
assess the benefits of using the $\SPN$ model over the simpler diffusion model, the author also implemented (still on
top of the neutronics framework) a discontinuous Galerkin discretization of the \textit{discrete ordinates}
approximation of the transport equation (the $\SN$ method, introduced in \Sref{sec:1-SN}), which, theoretically as $N\to\infty$, converges to the true solution of
the NTE (unlike the $\SPN$ approximation in multidimensional, multiregion domains). This will be addressed in Sec.
\alert{ref}.

\myparagraph{3D coupled neutron-physical finite element code based on the multigroup diffusion approximation}
Unfortunately, the 3D version of Hermes has not yet been released and thus the main utility of the neutronics framework
introduced in previous paragraph is in testing new ideas on hp-adaptivity for neutron transport problems. For the
purposes of the research project underlying this thesis (see the footnote \ref{ftn:TACR} on pg. \pageref{ftn:TACR}), a
3D finite element framework was needed. The final choice has been the FEniCS/Dolfin framework (\cite{dolfin1, dolfin2}).
The final code, that the author developed within this framework, is described in Chap. \alert(ref) and demonstrated on
several industrial and custom benchmarks. It allows parallel assembly of the multigroup neutron diffusion problem and
solution of the obtained algebraic problem using the well-established PETSc/SLEPc solvers (\cite{petsc1, slepc1})
wrapped by FEniCS.
This can be repeated in a feedback loop (until a steady state is found), in which the computed neutron flux directly influences
thermal/hydraulic properties of the core, the change of which in turn leads to a change of coefficients in the diffusion
equations. On top of that loop, another loop representing fuel burnup can be executed. At this point, the author would
like to acknowledge the work of his colleagues -- R. Ku{\v z}el (the coordinator of the whole effort and also the author
of an experimental eigensolver module), J. Egermaier and H. Kopincov{\' a} (who implemented the thermal/hydraulics
module) and Z. Vastl (who generated the meshes for the MOX-UO$_2$ and VR-1 benchmarks).

\myparagraph{Well-posedness of the $\SPN$ equations}
While there are many papers where the finite element method has been \textit{used} to find the weak solution of
both diffusion and $\SPN$ equations (e.g., \cite{Ragusa1, Hermes-nuclear, Ragusa2}), the question of well-posedness of
the corresponding variational formulation is neither addressed in these papers, nor do they provide references that
answer it. In the case of multigroup diffusion, well-posedness has been proved in \cite{Bourhrara1}. However, the author
couldn't find any paper dealing with the well-posedness of the $\SPN$ equations. It is therefore proved in Sec.
\alert{ref}.

\myparagraph{New method for solving large eigenvalue problems}
Determination of reactor steady state requires finding the dominant eigenpair of a generalized eigenvalue problem. This
has been traditionally performed by the simple power method, accelerated by eigenvalue shifting or Chebyshev
combination of previous iterates. Recently, Krylov subspace methods have been applied (\cite{warsa}, \cite{Subramanian},
or basically any reference from the ``Nuclear Engineering'' section at
\url{http://www.grycap.upv.es/slepc/material/appli.htm}). While the SLEPc library also serves as the main eigenvalue
solver of our neutron-physical code, there is another eigensolver module that is based on a new method
proposed by author's colleagues P. Van{\v e}k and R. Ku{\v z}el (who implemented the method on graphical accelerator
cards). Even though the implementation is not yet competitive with the well-optimized
solvers from the SLEPc library (in terms of execution speed), the method itself has some interesting properties
connected with solving large-scale eigenvalue problems. The paper describing this development has not yet been 
published, so it is included here as Chap. \alert{ref}. This last chapter also includes numerical experiments on simpler
problems, which confirm the theory and have not been included in the original paper.



  \chapter{Review of neutron transport methods}\label{chap:nte-review}

\nomenclature[s]{$f \approx g$}{$f$ is approximated by $g$}
\nomenclature[s]{$f \equiv g$}{$f$ is by definition equivalent to $g$}
\nomenclature[a]{$\R[n]$}{an n-dimensional Euclidean vector space}

The steady state neutron transport equation is a mathematical representation of balance between neutron gains and losses
within a given bounded macroscopic domain $\VV\subset\R[3]$. Let us consider the equation in its integro-differential
form with given neutron source function $Q$:
\begin{equation}\label{eq1}
  \begin{multlined}
    \Biggl[
      \bomega\cdot\grad + \sigma(\br,E)
    \Biggr]
    \psi(\br,\bomega,E) =\\
    = \intE[']{\Emin}{\Emax}{\intA[']{\kappa(\br,\bomega\sla\bomega',E\sla E')\psi(\br,\bomega',E',t)}}  + \src(\br,\bomega,E).
  \end{multlined}  
\end{equation}
It's solution, the \textit{angular neutron flux density} $\psi$ -- is a function of the following independent variables
that define the neutron phase space:
\begin{itemize}
 	\item $\br = (x,y,z)$
 	\nomenclature[A]{\br}{position vector}
 	\nomenclature[U]{x,y,z}{components of vectors in Cartesian coordinate system} represents the spatial distribution of
 	 neutrons,
 	\item $\bomega$\nomenclature[g]{$\bomega$}{unit vector of neutrons flow direction} represents the angular
 	distribution of neutrons on a unit sphere $\Sphere$\nomenclature[A]{\Sphere}{unit sphere ($\{x\in\R[3]: \norm{x} =
 	1\}$)} , i.e. their flow direction
 	\item $E\in [E_{\text{min}},E_{\text{max}}]$\nomenclature[g]{$E$}{Energy of neutrons\nomunit{eV}} is the kinetic
 	energy of neutrons.
\end{itemize}
Note that since the direction vectors are confined to the sphere: $\bomega\in\Sphere$, we can express the three
Cartesian components of $\bomega$ by only two spherical coordinates $\polar$ and
$\azimuthal$\nomenclature[g]{$\polar$}{polar angle}\nomenclature[g]{$\azimuthal$}{azimuthal angle}:
\begin{equation*}
	\bomega = \left[\begin{array}{c}
		\Omega_x \\
		\Omega_y \\
		\Omega_z
	\end{array}\right] = \left[\begin{array}{c}
		\sint\cosp \\
		\sint\sinp \\
		\cost
	\end{array}\right].
\end{equation*}

Function $\sigma$ represents reactions that result in a loss of neutron, while $\kappa$ represents reactions that
introduce neutrons into direction $\bomega$ and energy $E$ by, e.g., scattering from direction $\bomega'$ or fission. 
They are given by material composition of the domain and we will return to their more detailed description, as
well as to boundary conditions for eq. \eqref{eq1}, in section \ref{chap:NTE}. We just note at this moment that in our macroscopic
description, we always consider beams of neutrons with the same given average properties, even though we usually refer
to them as single neutrons with particular position, energy or direction. Similarly, we also omit the ``density''
specification of repeatedly used quantities, e.g. we will henceforth call $\psi$ just \textit{angular neutron flux}.

Later, we will also consider equation \eqref{eq1} without the source term $Q$ and with homogeneous boundary conditions.
Equation \eqref{eq1} then becomes a homogeneous integro-differential equation and finding its non-zero
solution will require solution of an eigenvalue problem. This reflects the fact that without external sources of neutrons, the exact balance between neutron gains and
losses within the bounded domain is possible only under precise physical conditions (the departure from which is
characterized by the corresponding eigenvalue). Ensuring such conditions is the goal of the core design and fuel
reloading optimization and this eigenvalue problem thus plays a vital role in the associated calculations. In
\Cref{chap:EVC}, we will introduce a new method for solving such problems that is particularly suited for solving
large-scale eigenvalue problems. As practically all numerical methods for solving eigenvalue problems involve the
solution of fixed-source problems of type \eqref{eq1}, we will use eq. \eqref{eq1} as our starting point.

Because of the relatively high number of independent variables appearing in the solution of eq. \eqref{eq1}
($x,y,z\vartheta,\varphi,E$), its practical use for calculations in larger spatial domains has always been a difficult
task. Various simplifications have therefore been explored and we review the most successful ones and their most
important properties in the following sections.

\section{Approximation of energetic dependence}

The continuous dependence on energy, $\psi = \psi(\cdot, E)$, is typically resolved by the so called \textit{multigroup
approximation}. In this approximation, eq. \eqref{eq1} is transformed into a finite system of integro-differential
equations, each governing the flux of neutrons with energies within a particular range (in this context called
\textit{group}):
\begin{equation}\label{eq:psi-MG}
  \psi^g(\br, \bomega) = \intEg{g-1/2}{g+1/2}{\psi(\br, \bomega, E)},\quad g = 1, 2,\ldots G.
\end{equation} 
These equations are coupled through reactions caused by neutrons with different energies, as given by the
first term on the right hand side of eq. \eqref{eq1}. Although the multigroup system of neutron transport equations has a
relatively simple form, finding an optimal grouping of energies and determining the associated group-averaged
coefficients is not an easy task in most practical applications because of the highly complicated energetic dependence
of nuclear processes. Even though an alternative to the finite-volume like approximation \eqref{eq:psi-MG} has been
proposed recently (\cite{Douglass}) -- using Galerkin projection of angular flux onto a space of 
functions supported over subregions of the energy range (a finite-element like approach) --  the multigroup
approximation still remains the most universally used approach to simplify the energetic dependence (see, e.g.,
\cite[Chap.~5]{Cacuci1} or \cite{Cho1}). In the remainder of our introduction of the neutron transport equation, we
will focus on the approximation of neutron flux in a single group (index of which will be omitted), described by the 
corresponding within-group equation in which contributions from other groups are encapsulated in the source term $\src$.
We will return to the multigroup approximation in \alert{ref}.

\section{Approximation of angular dependence}

Up to this day, a rather big number of methods have been proposed for approximating the angular dependence of neutron
flux. Many of them are still used and actively developed today as their characteristics make them more suitable for one
application area than other methods, which are preferred in different areas\footnote{of course, this does not exclude
their use in the other areas}.

\subsection{Lattice calculations}
As a first example, we consider the class of
methods originally derived from the equivalent integral form of the NTE (\alert{ref}). Typical representatives of this
class are the method of collision probabilities or the method of characteristics (see e.g.
\cite{Cho2,Wu1,Hursin1,Petkov1,Sanchez1}). As the integral form of the NTE represents the global neutron balance over
the domain (see \alert{ref}), the corresponding algebraic systems (obtained after spatial discretization) are full and
their solution demanding on computer resources. On the other hand, these methods quite naturally handle complex
geometries. Taking into consideration smaller geometric features of the domain, we are effectively coming from a
macroscopic scale to a mesoscopic range where the neutrons direction of motion as well as their kinetic energy become
more significant. High degree of spatial coupling and the requirement of fine resolution of angular and energetic
dependence does not make these methods suitable for whole-core reactor calculations.
However, to make the repeated whole-core simulations efficient and sufficiently accurate, these small-scale,
high-fidelity calculations\footnote{called \textit{lattice calculations} as they are typically performed on a single
representative subdomain of the core (one or several neighboring assemblies of fuel pins, or the fuel pin itself)
with reflective boundary conditions, simulating an infinite lattice of such subdomains} are indispensable for
generating appropriately averaged coefficients for the computationally more feasible larger scale.
This \textit{spatial homogenization} and \textit{energy group condensation}, as these averaging procedures are
traditionally called in nuclear engineering field, are employed by many existing whole-core simulators (see e.g.
\cite[Chap. 17]{Reuss1} or the review in first two sections of \cite{Sanchez7}). To simulate a long-term nuclear reactor
operation, it is furthermore neccessary to perform these procedures under varying physical conditions of the core and
generate many sets of averaged coefficients corresponding to these conditions. The code system described in
\alert{chapter} expects these coefficient sets on input, i.e., it is not designed for lattice calculations.
\begin{center}
***
\end{center}
More suitable for whole-core calculations (and hence for the purposes of the project underlying the authors
Ph.D. research) are methods derived from the integro-differential
version of the NTE, eq.
\eqref{eq1}, which lead to sparse algebraic systems. The most successful and nowadays best established are the method
of discrete ordinates ($\SN$) and the method of spherical harmonics ($\PN$).
Both these methods arise from applying a classical well-known approach for constructing finite numerical
approximations of PDEs  in the
angular domain. Although in the final code, we ultimately use the lowest order approximation that can be obtained
equivalently from both approaches -- the well-known neutron diffusion equation -- we will briefly introduce the general
$\SN$ and $\PN$ methods and expose their most important properties in the next two subsections. These properties are
well known, but their origin in mathematical structure of the equations is often overlooked in papers (rare
exceptions will be cited below). Also note that both the $\SN$ and $\PN$ methods have an importance for the thesis -- 
the author implemented the former in order to assess accuracy of the lower-order methods (\alert{section}) and developed
an alternative to the $\PN$ method that will be described in \alert{section}.

\subsection{The $\SN$ method}\label{sec:1-SN}
The $\SN$ method uses the collocation approach in which a set of directions (\textit{ordinates})
$\{\bomega_n\}_{n=1}^{M(N)}$ is chosen and the solution is approximated as:
\begin{equation}\label{eq:sn_approx} 
	\psi(\br, \bomega) \approx \pw{\psi(\br,\bomega_n)}{ when $\bomega = \bomega_n$}{0}{otherwise},\quad
	\text{ for each } n = 1,2,\ldots,M(N) 
\end{equation}
Equation \eqref{eq1} is then evaluated at these $M(N))$ isolated directions. Note
that for the traditional direction sets, we have $M = N(N+2)$ if the given problem does not possess any symmetries; the
method of discrete ordinates using such a number of directions is traditionally refered to as the method of discrete
ordinates of order $N$, shortly $\SN$. We will keep this notation, but  write just $M$ for the number of directions
instead of $M(N)$.

To evaluate the integral term on the right hand side of the equation, the set of directions is accompanied by a
corresponding set of weights $\{w_n\}_{n=1}^M$, together defining a quadrature of the sphere $\Sphere.$ The requirement
of accurate evaluation of the integral term of eq.
\eqref{eq1} as well as an accurate integration of the angular flux function over the sphere (which defines the
physically important quantity \textit{scalar flux}, \alert{ref}) constitutes the main guideline for the choice of
directions and weights. We will return to this matter later in \alert{ref}; for now it suffices to say that for
three-dimensional problems without any symmetries, \mbox{$M = \left|\{\bomega_n, w_n\}\right| = \Oh(N^2)$} (with the
value of $M$ for typically used quadrature sets stated above).

To write  the final result of the $\SN$ approximation, let us first denote for a set $s = \{c_k\}_{k=1}^n$ the column
vector with entries $c_1,c_2,\ldots,c_n$ as $\col(s)$ and similarly the diagonal matrix whose diagonal is
given by elements of $s$ as $\diag(s)$. Then, for
 $$\psi_n(\br) \equiv \psi(\br,\bomega_n), \quad q_n(\br) \equiv
q(\br,\bomega_n),$$ we will denote the vector of angular fluxes $\Psi(\br) = \col(\{\psi_n(\br)\})$ and the vector of
angular sources $\mat{Q}^{SN}(\br) = \col(\{q_n(\br)\})$.
The $\SN$ approximation consists of the following set of $M$ PDEs in spatial domain \footnote{Differentiation and
integration of vector functions (such as the term $\pd{\Psi(\br)}{x}$ in eq. \eqref{eq:sn1}) will be understood
component-wise throughout this thesis.}:
\begin{equation}\label{eq:sn1} 
\mat{A}^{SN}_x\pd{\Psi(\br)}{x} + \mat{A}^{SN}_y\pd{\Psi(\br)}{y} +
\mat{A}^{SN}_z\pd{\Psi(\br)}{z} + (\sigma(\br)\mat{I} - \mat{K}^{SN}(\br))\Psi(\br) = \mat{Q}^{SN}(\br),
\end{equation}
where $\mat{I}$ is the $M\times M$ identity matrix and
$$
	\mat{A}^{SN}_x = \diag(\{\Omega_{n,x}\}),\ \mat{A}^{SN}_y = \diag(\{\Omega_{n,y}\}),\ \mat{A}^{SN}_z =
	\diag(\{\Omega_{n,z}\}).
$$
Equation
\eqref{eq:sn1} represents a system of advection-reaction equations with constant advection field given by the angular 
matrices $\mat{A}^{SN}_x, \mat{A}^{SN}_y, \mat{A}^{SN}_z$. The system is weakly coupled through the absolute term
$\mat{K}^{SN}(\br)$, which is the quadrature representation of the original integral term, but not through the
differential terms. This fact and the relative smallness of the coupling term is utilized in a classical solution
technique for the $\SN$ approximation -- the \textit{source iteration} -- in which the system \eqref{eq:sn1} is fully
decoupled. Each equation is solved separately using any method suitable for an advection-reaction PDE, using $\psi_n$
from previous iteration to evaluate $\mat{K}^{SN}\Psi$ (both Jacobi and Gauss-Seidel schemes may be used to update
$\Psi$ during the iteration process).

\subsubsection{Ray effects}
A major drawback of the $\SN$ angular approximation comes from the fact that it assumes the radiation to propagate only
in a discrete set of directions. Only in the limiting case of infinite number of directions will the whole phase space
be covered. Otherwise, there will remain under-treated regions, while other regions will receive more radiation in order
to satisfy the global balance. This will lead to spatial oscillations of the scalar flux (obtained by adding up the
contributions from the individual directions) known as the \textit{ray effects}. This problem may be seen to be the
consequence of weak coupling between the unknowns in the $\SN$ system, which is however one of its biggest strenghts
from the computational point of view\footnote{This also exhibits the class of problems suffering the most from ray
effects -- namely problems without fission and scattering that lack the coupling term $\mat{K}^SN$} .
It is therefore very difficult to reduce the ray effects while keeping the advantages of the $\SN$ approximation (for an
overview and some heuristic observations, see \cite{Li1}). 

The root cause of ray effects in the $\SN$ approximation is the loss of an inherent property of the radiative transfer
equation \eqref{eq1} when transforming it into the system \eqref{eq:sn1}. This property is the rotational invariance: If
we write eq. \eqref{eq1} in operator form as 
\begin{equation}\label{eq1op}
	\mathcal{L}\psi = \mathcal{K}\psi + \mathcal{Q}
\end{equation} 
(the precise mathematical setting for the operators will be given in \alert{ref}), then it is not
difficult to see that 
$$
	\mathcal{R}(\mathcal{L} - \mathcal{K}) = (\mathcal{L} - \mathcal{K})\mathcal{R}
$$ 
for an operator $\mathcal{R}: f(\br,\bomega) \mapsto f(R^T\br,R^T\bomega)$ corresponding to any rotation of
coordinate system $R$, provided that the coefficient functions $\sigma$ and $\kappa$ are also invariant under the action
of $\mathcal{R}$. This will be the case, e.g., if both $\br$ and $R^T\br$ lie in an isotropic\footnote{Note that this
does not mean that the actual process of scattering is isotropic, see \alert{ref}.}, homogeneous region -- a typical
situation in multi-zone reactor calculations (see, e.g., \alert{Zweifel, Hauck}). If we define 
$$
	\mathcal{P}^{SN}: \psi(\br,\bomega) \mapsto \Psi(\br) 
$$
by eq. \eqref{eq:sn_approx} and again formally express eq. \eqref{eq:sn1} in the operator form:
$$
	\bigl(\mathcal{L}^{SN} - \mathcal{K}^{SN}\bigr)\mathcal{P}^{SN}\psi = \mathcal{Q}^{SN},
$$
(with the same assumption on $\sigma$ and $\kappa$), its clear that now
$$
	\mathcal{R}\bigl(\mathcal{L}^{SN} - \mathcal{K}^{SN}\bigr)\mathcal{P}^{SN} \neq \bigl(\mathcal{L}^{SN} -
	\mathcal{K}^{SN}\bigr)\mathcal{P}^{SN}\mathcal{R},
$$	
because $\mathcal{P}^{SN}\mathcal{R}\psi(\br,\bomega) = 0$ if $R^T\bomega \not\in \{\bomega_n\}_{n=1}^{M}$. As a
consequence, the ray effects cannot be completely avoided in the standard $\SN$ framework. This brings us to the second
 approach used for approximating angular dependence in the integro-differential NTE that naturally overcomes
this issue.

\subsection{The $\PN$ method}

Instead of the collocation method used by the $\SN$ approximation, the $\PN$ method uses the Galerkin method in
angular domain. That is, the angularly dependent quantities in \eqref{eq1} are expanded into infinite series of properly
chosen functions that span a complete basis on the unit sphere, the equation is multiplied by each member of the basis in turn and integrated over the sphere. The properties of the basis functions are then used to derive
equations for the expansion coefficients. Only a finite number of the expansion terms is considered to allow practical
computation. Usually, the expansion is truncated at length $N$ by setting all expansion coefficients with higher index
to 0 (although there exist alternative closure methods that may have favorable properties in certain situations, see
e.g. \cite{Frank0}). Then we speak of the $\PN$ approximation:
\begin{equation}\label{eq1.1}
  \psi(\br,\bomega) \approx \sum_{k=1}^{\widetilde M} \phi_k(\br) f_k(\bomega).
\end{equation}
The set of spherical basis functions that were used in the original $\PN$ method are the
\textit{spherical harmonic functions}, which simplify the algebraic manipulations needed to arrive at the conditions 
for coefficients $\psi_k$ (called \textit{angular moments}) and have other favorable properties (some will be mentioned
in \alert(ref)). In one dimension, the spherical harmonic functions reduce to Legendre polynomials and $\widetilde M = N$. For general three-dimensional
problems, there are $2n + 1$ linearly independent spherical harmonics for each degree $n$ and
$$
	\widetilde M = \sum_{n=0}^{N} 2n + 1 = (N+1)^2.
$$
We expand the approximation \eqref{eq1.1} as
\begin{equation}\label{eq:pn_approx}
	\psi(\br,\bomega) \approx \suma[n]{0}{N}\suma[m]{-n}{n}\angmom{n}{m}(\br)\Y{n}{m}(\bomega)
\end{equation}
where $\Y{n}{m}(\bomega)$ is the spherical harmonic function of degree $n$ and order $m$. Therefore, in \eqref{eq1.1},
we consider the single index $k$ ($1 \leq k \leq \widetilde M$) that covers all the combinations of $n$ and $m$ ($0 \leq n
\leq N$, $-n\leq m \leq n$) appearing in \eqref{eq:pn_approx}. We finally arrive at a system of $\widetilde M$ partial
differential equations in spatial domain which is of comparable size as the system of $\SN$ equations and has the following form:
\begin{equation}\label{eq:pn1}
	\mat{A}^{PN}_x\pd{\Phi}{x} + \mat{A}^{PN}_y\pd{\Phi}{y} + \mat{A}^{PN}_z\pd{\Phi}{z} + (\sigma\mat{I} -
	\mat{K}^{PN})\Phi = \mat{Q}^{PN},
\end{equation}
where $\Phi = \col(\{\phi_k\})$ and $\mat{Q}_{PN} = \col(\{q_k\})$ are, respectively, the vector of angular flux
moments and analogously defined angular source moments. The system has the same form as that for
the $\SN$ approximation, but the angular matrices:
$$
	\left[A^{PN}_s\right]_{k,l} = \intA{\Omega_s \Y{k}{}(\bomega)\Yc{l}{}(\bomega)},\quad s\in\{x,y,z\},\ 
	1 \leq k,l \leq \widetilde M
$$
(overbar denotes complex conjugation) are no longer diagonal, but rather banded with each moment being coupled to at
most six other moments (see, e.g., \cite{Sanchez7}).

\subsubsection{Rotational invariance of the $\PN$ equations}
The increased coupling between the unknowns in the $\PN$ system is the price for the attractive property of spherical
harmonics that prevents ray effects appearing in $\SN$ solutions. It is well known (e.g, \cite[Chap.3]{Sansone}) that
the mutually orthogonal spherical harmonic functions of given degree $n$ generate a rotationally invariant subspace of
the Hilbert space of square-integrable functions on the sphere $\Lp{2}(\Sphere)$, which we denote by $\Lambda_n$:
\begin{equation}
	\label{eq:subspace}
    	\Lambda_n = \mathrm{Span}\bigl\{\Y{n}{m}; -n \leq m \leq n\bigr\},
\end{equation}
This means that $\mathcal{R}(\Lambda_n) \subset \Lambda_n$ for the rotation transformation $\mathcal{R}$ (cf. the
paragraph on ray effects in \Sref{sec:1-SN}). Moreover, $\Lambda_n$ is the eigenspace associated with the $n$-th 
eigenvalue of the Laplace operator on $\Sphere$, $\lambda_n = -n(n+1)$. Being eigenspaces of a self-adjoint operator
corresponding to different eigenvalues, $\Lambda_n$ for $n=0,1,\ldots$ are mutually orthogonal and  
$\Lp{2}(\Sphere) = \bigoplus_{n=0}^{\infty}\Lambda_n$\nomenclature[S]{$\bigoplus$}{direct sum of spaces}.
Hence, we can write the $\PN$ projection onto the finite-dimensional subspace $\Lp{2}_N\subset\Lp{2}(\Sphere)$, eq.
\eqref{eq:pn_approx}, as
\begin{equation}\label{eq:pn_approx2}
	\psi(\br,\bomega) \approx \mathcal{P}^{PN}\psi(\br,\bomega) = \suma[n]{0}{N}\angmom{n}{}(\br,\bomega),
\end{equation}
where 
$$
	\angmom{n}{}(\br,\bomega) = \suma[m]{-n}{n}\angmom{n}{m}(\br)\Y{n}{m}(\bomega) \in \Lambda_n
$$
with uniquely determined coefficients $\angmom{n}{m}$. Using linearity of the rotation operator and rotational
invariance of each $\Lambda_n$, it follows that 
$$
	\mathcal{R}\mathcal{P}^{PN} = \mathcal{P}^{PN}\mathcal{R}
$$
and, since the $\PN$ system \eqref{eq:pn1} is actually obtained by applying $\mathcal{P}^{PN}$ to the left and right
hand side of eq. \eqref{eq1} and using linear independence of the spherical harmonic functions. 

\subsubsection{Drawbacks of the $\PN$ approximation}
Using the results of the previous paragraph and well-known results from the theory of Hilbert spaces, we can see that
the sum \eqref{eq:pn_approx2} (or \eqref{eq:pn_approx}) converges in the $\Lp{2}(\Sphere)$ norm to the true solution of
eq. \eqref{eq1} as $N\to\infty$. However, the convergence may be very slow if the true solution to the NTE is not
sufficiently regular in the angular variable. In particular, pointwise convergence is hindered in the neighborhood of
phase space points where the solution of the NTE has jump discontinuity in $\bomega$ (which may occur for example when a
narrow beam of neutrons is freely streaming through a non-interacting medium, but also in a more typical case of domains
with multiple regions with different materials, bounded by piecewise polygonal boundary; see also \alert{discussion on
regularity properties of NTE}) and spurious oscillations are introduced to the approximate solution at these points.
These oscillations spread over the whole angular domain and slow down the norm-wise convergence. This is a well-known
property of Fourier series (which the expansion \eqref{eq:pn_approx2} generalize) known as \textit{Gibbs phenomenon}.
Moreover, these oscillations do not vanish as more terms in the series are retained. However, there are several ways of
circumventing the Gibbs phenomenon. For example, when considering \eqref{eq:pn_approx2} as a means of deriving the $\PN$
system, we may note that using a finite expansion obtained by truncating \eqref{eq:pn_approx2} at $n=N$ is not the only
way of obtaining a closed system of equations -- different closures are possible as already discussed above. This fact
has been utilized in \cite{McClarren3} where the expansion has been adjusted to mitigate the oscillations by controlling
angular gradients\footnote{Note that the expansion \eqref{eq:pn_approx2} represents the best $\Lp{2}(\Sphere)$
approximation of $\psi$ by spherical polynomials up to a given degree, but absence of angular gradients in the
$\Lp{2}(\Sphere)$ norm permits arbitrary oscillations.}. For other similar approaches in the context of general spectral
methods, see e.g. \cite{Tanner}.

As shown in \cite{McClarren4}, there is also another issue connected with time-dependent $\PN$ approximation that must
be kept in mind particularly when solving coupled problems. This issue is inherent in the structure of the $\PN$ system
and cannot be completely removed without losing some of its attractive properties. Namely, the authors proved that
without sources and reactions, the linear hyperbolic character of eq. \eqref{eq:pn1} (with an additional time derivative
term) together with rotational invariance allows negative solutions for positive, isotropic data in two or three
dimensions. To prevent negative solutions, one could either give-up linearity (e.g. by using a non-linear closure in a
similar way as described above), rotational invariance (thus introducing ray-effects into the solutions) or
hyperbolicity (thus changing the speed at which radiation propagates throughout the domain) -- none of which is a
generally satisfactory remedy. The authors also demonstrated that negative solutions can appear even in heterogeneous
domains containing regions with reactions or sources.

\section{Approximation of spatial dependence in $\SN$ and $\PN$ methods}
As we have seen in the previous subsections, both the $\SN$ and $\PN$ approximations lead to a system of linear
hyperbolic PDE's in spatial variables. The final approximation step typically consists of laying out a mesh over the
spatial domain and using finite difference (FD), finite volume (FV) or finite element (FE) methods to discretize the
PDE's. In the case of the $\SN$ approximation, the approach traditionally favored by the nuclear engineering community
uses the source iteration technique to decouple the system into single-direction equations, each of which is then solved
by a \textit{transport sweep} from inflow to outflow boundaries of mesh cells. The finite volume method is used to link
the mesh cells through cell-averaged and interface unknowns. Without fission and scattering (i.e., the integral term
$\mathcal{K}$ in \eqref{eq1op}) and reflective boundaries (\alert(ref)), the system is already decoupled and only one
sweep for each direction is sufficient to determine $\mathcal{L}^{-1}$ and hence the solution. In the other case, source
iteration is required and when the integral term is dominant (like in the case of whole-core reactor calculations) some
acceleration is usually required (see the discussion in Sec. \ref{:}).

Similarly to the explicit schemes for usual time-dependent advection-reaction problems, this direction-sweeping scheme
requires careful choice of numerical approximation of interface unknowns to ensure stability (restricting and
intertwinning the angular and spatial resolution). An alternative, attractive in particular when an unstructured
spatial mesh is used (and even more in three dimensions), is the implicit solution of of the angularly discretized
system. This is also the method of choice for the $\PN$ system, where diagonalization of the angular matrices for each
differently oriented cell would be required for the sweeping procedure. However, the stability constraints do not
disappear completely -- in order not to introduce unwanted oscillations into the solution, a stable spatial
discretization must be used to obtain the algebraic system. In the case of the finite element method, for example, it is
well known that the approximation of the solution of an advection-reaction PDE by piecewise continuous functions (i.e.,
the continuous Galerkin method) is not stable and allows arbitrarily large derivatives of the solution in the direction
of flow (and hence the oscillations, see \alert{ref}).
Therefore, either the discontinuous Galerkin (DG) method or some version of stabilized continuous Galerkin method are
required for spatial discretization. We will describe the simplest DG method for the $\SN$ system in \alert{section};
(see e.g.
\cite{Meinkohn} for the application of the streamline-upwind Petrov Galerkin method).

In any case, the final system of linear algebraic equations is generally sparse but -- given the complicated geometry
and material arrangement of realistic problems -- very large. It is usually computationally infeasible to resolve all
local features of the solution by a uniform mesh and some sort of adaptivity is employed. In some cases (typically in
engineering application like the simulation of heterogeneous nuclear reactor cores), the long-term experience may be
used to create the mesh by hand using well-established geometry and mesh generation software like the commercial CUBIT
(\cite{CUBIT}) or the open-source GMSH (\cite{GMSH}). When the a-priori knowledge of important solution feautres is not
available, some sort of automatic adaptivity needs to be employed, which we will discuss in the following section.
Nevertheless, both approaches lead to very large systems (easily of the order of $10^8$ even for a crude energetic
discretization), often ill-conditioned as a result of highly irregular material properties or anisotropic meshes
(particularly when automatic mesh refinement is employed). Using even the advanced sparse direct solvers like
UMFPACK or MUMPS (\cite{UMFPACK,MUMPS}) to solve such systems is not practical.
Moreover, it is intuitively obvious and easily proved (see e.g. \cite{Arioli}) that 

Moreover, there are cases (typically
in engineering applications like the simulation of heterogeneous nuclear reactor cores) where it is common to create
initial spatial mesh by some specialized CAD and mesh-generation system based on the long-time operating experience. ,
The need to resolve a complicated dependence on 6 independent variables may easily result in systems with

Robust solvers that Also, because of the possibly highly heterogeneous domains with large jumps in material coefficients
between subdomains, conditioned

\subsection{Adaptivity}
In real applications, automatic adaptivity of the discrete phase space is usually needed to obtain sufficiently accurate
results sufficiently fast. Except for the scheme described in the Ph.D. thesis of H. Park \cite{Park} -- a rare case of coupled angular and spatial adaptivity -- all literature
available to the author describes schemes where angular and spatial adaptivity is performed separately (and there is
none about adaptivity with respect to the energetic variable). In the context of the $\PN$ method, there are examples
where the order of the spherical harmonic expansion is varied throughout the spatial domain (see e.g. \cite{Ackroyd2})
based on material properties and physical reasoning. Properties of spherical wavelets have been used in \cite{Buchan} to
drive automatic selection of the order of the expansion (in this case with respect to the wavelet interpolation basis of
$\Lp[2](\Sphere)$ instead of the spherical harmonics basis) based on the increasing size of expansion coefficients
corresponding to wavelets supported over underresolved areas of $\Sphere$. Angular adaptivity in the $\SN$ approximation
(i.e., adaptive control of the number of discrete ordinates in local areas on the sphere) is described in
\cite{Jarrell}.

More widespread use has found the adaptivity in spatial domain, using techniques developed for finite element
approximations of general hyperbolic systems. They are usually used in DG $\SN$ methods, see e.g.
\cite{Fournier,Duo,ragusa2010two} for adaptivity based on a-posteriori estimation of global $\Lp[2]$ norm of solution
error or \cite{LathouwersGoal, Wang2} for a method based on goal oriented adaptivity. Spatial adaptivity for the
standard $\PN$ approximation is not used as widely, probably because the approximation itself is not so widely used for
larger scale calculations. However, there are various reformulations of the $\PN$ system as a system of second-order
PDEs for which many usual a-posteriori estimates for
elliptic PDEs have been successfully applied. These formulations will be introduced in the
following section.

The discontinuous Galerkin framework used by majority of $\SN$ methods doesn't impose any constraints on the solution
continuity accross elements and allows the solution to be represented by completely different functions on each element.
As such, it is well-suited for implementing both the mesh refinement and polynomial order variation procedure, paving a
way for hp-adaptive FE solution.
Nevertheless, all the references above are employing h-adaptivity where the mesh is refined with fixed polynomial
approximation degree $p$. Prevalence of h-adaptivity and use of $p=1$ (linear) finite element spaces is caused partly
by the difficulty of implementation of an hp-adaptive FE code itself, partly by the fear of the well known limited
regularity of the exact solution of \eqref{eq1} even for smooth input data\footnote{By the method of characteristics, we
can expect the angular fluxes $\angflux(\cdot,\bomega)$ to be differentiable in the direction of $\bomega$, but not in
any other direction. As shown in \cite{Johnson}, the scalar fluxes, i.e. integrals of $\angflux$ over $\Sphere$, belong
at most to $H^{3/2-\epsilon}(\VV)$ where $0 < \epsilon \ll 1$ and $H^{k}(\VV)$ denotes the usual Sobolev space of order
$k$.}.
However, similarly to the experience with hp-adaptive methods in different fields, the limitation of asymptotic
convergence rate (as $h\to 0$ where $h$ is the diameter of the largest element in the mesh) dictated by a-priori error
estimates involving solution regularity typically doesn't appear until very late in the mesh refinement process or at
all (\cite{wang2009convergence}). Hence, utilizing higher order approximations still makes sense to accelerate
pre-asymptotic convergence rate as much as possible (for an attempt to use hp-adaptivity with DG $\SN$ methods -- to
author's knowledge first of its kind, see \cite{FournierDGHP}).



% REGULARITY


  

\subsection{Second-order formulations}

There is, however, another possibility to make the solution of the transport equation more feasible. The set of
steady state $\PN[1]$ equations, implying that the neutron flux varies only linearly in angle, can be under some
additional physically justifiable assumptions recast (even in 3D) into a single elliptic equation\footnote{The same
holds also for the $\SN[2]$ equations.}. This is the familiar \textit{diffusion approximation} -- thanks to its simplicity and also the
efficiency of the numerical solution techniques available for this approximation, it has always served as a ``workhorse
computational method of nuclear reactor physics" \cite[p. 43]{Stacey1}. The model is indeed sufficiently accurate for
whole core calculations of contemporary reactors, assuming that the significant finer-scale neutron transport processes
have been resolved by some higher-fidelity NTE solvers applied in previous solution stages. The self-adjoint diffusion
equation can then be solved using e.g. the finite element method in conjunction with both powerful and theoretically
well-established conjugate gradient method with symmetric preconditioners like the modern algebraic multigrid. Solution
efficiency may be improved even further also by using adaptive mesh refinement based on highly developed a posteriori
error estimates for elliptic problems. Note that the self-adjoint property of the diffusion model can only be spoiled by
the multigroup energy discretization, where energy transfers in neutron collisions result in non-symmetric coupling of
the multigroup system -- this can be however easily prevented by moving the non-symmetric parts to the right-hand side
and solving the resulting system iteratively. The simplest of such iterative splittings (which is nevertheless used by
many existing neutron transport codes, often in conjunction with some acceleration technique) is the so-called
\textit{source iteration}, which will be briefly introduced in Sec. \ref{sec:MG}.

Although methods based on the diffusion approximation have been experimentally proven to be widely applicable for
nuclear reactor analyses, there are situations where this approximation is just too coarse and, as some recent reports
indicate \cite{Hejzlar1,Cho1}, these cases are likely to grow soon with the advent of new reactor and fuel designs. This
approximation, of course, can also be hardly expected to produce acceptable results for more general problems with
strong transport effects (leading to steep solution gradients or even discontinuities), such as those arising in the
radiation shielding studies. One possibility then is to treat the diffusion solves not as a means of obtaining the final
solution, but as preconditioners in an iteration involving a rigorous transport solution. Particularly popular became
such coupling between the diffusion calculation and discrete ordinates source iteration, which got the name
\textit{diffusion synthetic acceleration} (\cite{Alcouffe1}). Research in this field is still very active, focusing for
instance on diffusion preconditioning of Krylov subspace iterations involving discrete ordinates matrices (\cite[Chap.
1]{Azmy1}).

Another possibility is to look for other ways of transforming the first-order neutron transport equation into a set of
second-order ones, keeping the favorable mathematical and numerical properties of the diffusion equation and yet not
losing the ability to capture the most important transport effects. The simplest approach appears to be to generalize
the procedure used to obtain the diffusion equation from the zeroth and first moment equations of the $\PN[1]$ set.
Although this leads to an attractive system of weakly coupled diffusion-reaction equations in 1D, a complicated system
of strongly coupled equations with mixed second-order partial derivatives results in more dimensions (\cite{Capilla}).
Another, more general approach uses alternative formulations of the NTE itself as an integro-differential equation with
second-order spatial derivatives. Probably best known of such formulations is the even (or odd) parity form. Its
derivation begins by writing eq. \eqref{eq1} for $\bomega$ and $-\bomega$ and adding and subtracting the two resulting
equations. Two first-order equations are thus obtained for the even/odd parity fluxes  $\psi(\cdot,\bomega)_{\pm} =
\tfrac12\bigl[\psi(\cdot,\bomega) \pm \psi(\cdot,-\bomega)\bigr]$, coupled only through zero-order terms. The parity
forms of the NTE are then obtained by eliminating either $\psi_+$ or $\psi_-$ between these two equations\footnote{The
even parity equation appears to be used more often than the odd parity one; it is in fact the Euler equation for a
functional that was used in the earliest variational characterization of neutron transport (see \cite[Chap. 2]{Azmy1}
and references therein).}. One can now apply any angular discretization technique (such as the above mentioned $\PN$ or
$\SN$ methods) on the equations thus obtained to construct a self-adjoint system amenable to stable continuous FEM
solution with the powerful algebraic solution techniques mentioned above.

An improvement of this parity formulation is the \textit{self-adjoint angular-flux equation} (SAAF, see \cite{Morel1}).
It removes some of the deficiencies of the parity equations caused by the fact that the latter provide only one (even or
odd) component and not the whole angular flux. Upon angular discretization, the SAAF equation is also related to the
weighted least-squares approximation, and to the Galerkin least-squares method (a popular stabilization technique for
continuous finite element approximation of general first-order hyperbolic PDE's). This latter form has the advantage of
natural error estimator provided by the least-squares functional, which may be used to drive spatial (or even angular)
adaptivity.

\subsubsection{The simplified $\PN$ approximation}

Although the significant accuracy improvements of the second-order approximations with respect to the diffusion equation
secure their place in difficult transport calculations, their usage for the analysis of contemporary nuclear reactor
cores still seems to be very limited. This is despite the fact that the presence of materials with very different
neutron-physical properties within these cores%
% \footnote{such as those loaded with the MOX/UO2 fuel}
pronounce the transport effects. More preferred method for such calculations is the \textit{simplified $\PN$
approximation ($\SPN$)} with origins dating back to the early 1960's (see the Gelbard's seminal papers
\cite{Gelbard1,Gelbard2}). Its derivation was completely formal at the beginning -- amounting to a simple replacement of
differential operators $\textstyle\der{}{z}$ in the 1D $\PN$ system by their multidimensional counterparts $\nabla$ and
$\nabla\cdot$ and recasting those scalar unknowns operated upon by the latter as vector quantities. Despite this
mathematically weak derivation, the $\SPN$ solution was known to be equivalent to the exact one in the case of a
homogeneous medium and, comparing to either diffusion or full $\PN$ models, provided encouraging results both in terms
of accuracy and efficiency even in more realistic cases. The method was also very attractive thanks to the relative
simplicity of implementation, requiring only modification of existing multigroup diffusion codes instead of writing
completely new ones (as was the case with the other transport methods mentioned above).

After some time, however, the analysts were able to devise special transport problems for which the simple diffusion
approximation actually provided better results (see, e.g., \cite[p. 247]{Coppa1}). Validity of Gelbard's formal
derivation therefore became questioned and the $\SPN$ equations have not been seriously considered as a robust enough
improvement of the diffusion model for some time. This has changed in the 1990's when the asymptotic and variational
analyses \cite{Larsen1,Brantley1,Pomraning1} theoretically justified the method and also rigorously determined the range
of validity of the approximation. Although it turned out that this range is not significantly larger than that of the
diffusion theory (\cite{Larsen1}), the $\SPN$ approximation has recently been shown to produce more correct results than
the diffusion model under these conditions and is gaining popularity again
\cite{Frank1,McClarren1,Ragusa1,Larsen3,Kirschenmann1}. Moreover, even though the $\SPN$ solution does not tend to the
exact solution of the NTE as $N\to\infty$ in general, there are several cases in which it is equivalent to the
convergent $\PN$ expansion (see some recent papers like \cite{Coppa2,McClarren2,Larsen2}) and further research of the
$\SPN$ model and its connections to the NTE appears to be an interesting topic.
	
%  \section{Organization of this report}

The steady-state Boltzmann transport equation with a fixed source term will be stated in Chapter 2 in both continuous 
energy and multigroup form. Following \cite{DautrayLions}, we will review conditions for the existence of a unique 
non-negative solution of the equation.

Main contribution of this preliminary dissertation report to the current state of neutron transport knowledge is 
contained in Chapter 3 where the \textit{Maxwell-Cartesian surface spherical harmonic functions} are used to derive an
alternative to the classical set of $\PN$ equations (the resulting set is called here ``$\MCPN$"). Despite its
favorable properties with respect to the traditional basis of tesseral spherical harmonics, these functions have not
been used in a systematic way for neutron transport problems before (the few related works will be dissected in Chap. 3
and main differences from the present work underlined). A very close relationship between the form of the $\MCPN$ 
equations and that of the 1D $\PN$ equations is then noted, possibly bringing new light into the question of validity
of the $\SPN$ approximation. The structure of the $\MCPN$ set also allows manipulations which could lead to other 
advantageous forms of the $\PN$ equations that were not thoroughly explored before -- one of them is outlined at the end
of the chapter.

Numerical methods for solving either the $\MCPN$ set or its derived forms (including the $\SPN$ equations) will be 
formulated and analyzed in the final dissertation thesis. Outline of this future work concludes this report.
		
%  \chapter{Mathematical model of neutron transport}\label{sec:bte}

\ifpdf
    \graphicspath{{2_BTE/figures/PNG/}{2_BTE/figures/PDF/}{2_BTE/figures/}}
\else
    \graphicspath{{2_BTE/figures/EPS/}{2_BTE/figures/}}
\fi

\section{Neutronic phase space} \label{sec:phase}
Continuing with the notation introduced in \ref{sec:bte-intro}, let us define the neutronic phase space $$
  X := \{(\br,\bomega,E):\ \br\in V\subset\R[3], \bomega\in \Sphere, E \in [\Emin,\Emax]\}
$$ together with its outflow and inflow boundary subsets, respectively:
$$
  \pX[\pm] := \bigl\{ (\br,\bomega,E) \in \pV \times \Sphere \times [E_m, E_M], \mbox{ s.t. } \bomega\cdot\bn(\br)
  \gtrless 0 \bigr\}
$$ We will assume that $V$ is a convex domain with piecewise smooth boundary $\pV$, oriented by its outward normal field
$\bn(\br)$. The direction vector $\bomega = \bv/v$ has the following Cartesian coordinates:
\begin{equation*}
	\bomega = \left[\begin{array}{c}
		\Omega_x \\
		\Omega_y \\
		\Omega_z
	\end{array}\right] = \left[\begin{array}{c}
		\sint\cosp \\
		\sint\sinp \\
		\cost
	\end{array}\right].
\end{equation*}
The product measure
\begin{equation}\label{eq:measure}
  \d{x} = \d{\mu(X)} = \d{\mu(V\times\Sphere\times [\Emin,\Emax])} = \d{\br}\d{\bomega}\d{E}
\end{equation}
is used when integrating over $X$, where $\d{\bomega}$ is defined as the solid angle subtended at the center of
 $\Sphere$ by the angular differential element:
$$
	\d{\bomega} = \sint \d{\polar}\d{\azimuthal}
$$
while the boundary measure
$$
\db = \abs{\bomega\cdot\bn}\d{\gamma}\d{\bomega}\d{E},\quad \d{\gamma} = \d{\mu(\pV)},
$$
is used when integrating over $\pX[\pm]$.

\section{Steady state neutron transport problem}\label{sec:definitions}

Equilibrium distribution of neutrons within the phase space, originating from a time-invariable volumetric source
$\src(\br,\bomega,E)$ within $V$, is given by the time-independent version of eq. \eqref{eq1}:
\begin{equation}\label{eq:bte1}
\begin{multlined}
  \bomega\cdot\nabla\psi(\br,\bomega,E) + \sigma_t(\br,E)\psi(\br,\bomega,E) =\\[.25em]
   = \intE[']{\Emin}{\Emax}{
      \intA[']{\kappa(\br,\bomega\cdot\bomega',E\sla E')\psi(\br,\bomega',E')}
    } + q(\br,\bomega,E).
 \end{multlined}
\end{equation}
together with specified distribution at the inflow boundary. The two most often used inflow boundary conditions are:
\begin{itemize}
	\item incoming angular neutron flux
	\begin{equation}\label{eq:bte2}
	  \psi\vert_{\pX[-]} = \psi_{\text{in}}
	\end{equation}
	($\psi_{\text{in}}\equiv 0$ corresponds to vacuum in $\R[3]\setminus \overline V$, which is a common
	 assumption in nuclear reactor modeling),
	
	\item specular reflection at boundary
	\begin{equation}\label{eq:bte3}
  	\psi(\br,\bomega,E) = \psi(\br, \bomega_R, E),\quad (\br,\bomega,E)\in \pX[-],\ \ \bomega_R = \bomega - 2 \bn 
  	(\bomega \cdot \bn) 
  \end{equation}
  (this condition is used to model planes of symmetry).
\end{itemize}
For more general boundary conditions used in conjunction with eq. \eqref{eq:bte1}, see e.g. \cite{Sanchez3}.

A physically plausible solution of BTE should be moreover non-negative throughout $V$ and continuous along any direction
$\bomega$, i.e. $\psi(\br + s\bomega,\bomega,E)$ is a continuous function of $s$ for any $\br$, $\bomega$, $E$. Note
that $\psi(\br + s\bomega',\bomega,E)$ \textsl{may} be discontinuous when $\bomega' \neq \bomega$.

The kernel of the integral operator on the right-hand side is commonly split as follows:
$$
  \kappa(\br,\bomega\cdot\bomega',E\sla E') = \sigma_s(\br,\bomega\cdot\bomega',E\sla E') +
  \frac{\chi(E)\nu(\br,E')\sigma_f(\br,E')}{4\pi}
$$ where $\sigma_s$ is called \textit{double-differential macroscopic cross-section for scattering} and $\sigma_f$ the
\textit{macroscopic cross-section for fission}. We note in the latter that fission is an isotropic process (thus
$\sigma_f$ is the average value over $\Sphere$), by which a neutron with energy $E'$ releases (on average) $\nu(\br,E')$
neutrons with new energy distribution $\chi(E)$. The new energy distribution is normalized so that $$
  \intE[']{\Emin}{\Emax}{\chi(E')\sigma_f(\br,E)} = \sigma_f(\br,E).
$$ New neutrons can also be released into the system as a result of inelastic scattering but this contribution is
usually neglected in applications of interest to us.

The \textit{total macroscopic cross-section}, $\sigma_t$, characterizing the probability that neutron with energy $E$
undergoes a collision of any type with nuclei at $\br$, can be decomposed as $$
  \sigma_t(\br,E) = \sigma_a(\br,E) + \sigma_s(\br,E) = \sigma_c(\br,E) + \sigma_f(\br,E) + \sigma_s(\br,E),
$$ where
\begin{itemize}
	\item $\sigma_a(\br, E)$ $\ldots$ absorption cross section,
	\item $\sigma_c(\br, E)$ $\ldots$ non-productive capture cross section,
	\item $\displaystyle\sigma_s(\br, E) = \intE[']{\Emin}{\Emax}{\intA[']{\sigma_s(\br,\bomega'\cdot\bomega, 
	E\shortrightarrow E')}}$.
\end{itemize}
A natural assumption is that the macroscopic cross-sections are bounded measurable functions
\footnote{under the product measure \eqref{eq:measure}, resp. under the product measure $\d{\mu(V\times\Sphere^2\times
 [\Emin,\Emax]^2)}$ for the double-differential cross-section}, piecewise continuous throughout $V$. 

From the solution of eq. \eqref{eq:bte1}, one can derive the following integral quantities
\begin{itemize}
  \item \textit{scalar neutron flux density}
  \begin{equation}%\label{eq:scalar_flux}
    \phi(\br, E) = \intA{\psi(\br,\bomega,E)},
  \end{equation}
  \item \textit{net neutron current density}
	\begin{equation}%\label{eq:bJ}
		\bJ(\br, E)	= \intA{\bomega\psi(\br,\bomega,E)},
	\end{equation}
\end{itemize}
so that integrating $\phi(\br, E)$ over a given volume or $\bJ(\br,E)\cdot\bn(\br)$ over a given surface, respectively,
gives the total number of neutrons with energy $E$ within that volume or crossing that surface in the direction of 
$\bn$, respectively. Also,
\begin{equation}\label{eq:rr}
  \sigma_x(\br, E) \phi(\br, E)
\end{equation}
represents the \textit{reaction rate} (per unit time) of given type. All these quantities produce directly observable 
effects and may be experimentally measured by various detector mechanisms, which is the reason why they are more 
important in practical calculations than the actual solution of the BTE (i.e. the angular neutron flux). Note that 
for the total volumetric scalar flux to be finite (as is physically expected), we should look for the solution $\psi$ 
of the BTE in the space $\Lp[1](X)$ with boundary values in $\Lp[1](\pX[\pm])$ (equipped with the respective measures
defined in Sec. \ref{sec:phase}), which is therefore usually taken as the natural function space setting 
(\cite{DautrayLions}; we will return to this topic in Sec. \ref{sec:solvability}).

Apart from the ``fixed source'' problem defined above, the other important problem in neutron transport requires the
determination of material composition (i.e. the values of $\sigma_x$) for a given reactor core geometry (or vice versa)
which ensures a steady neutron distribution (that means -- steady power generation) with no additional neutron sources
than the fission. This is called a ``criticality problem'' and entails finding a suitable parameter with which the
homogeneous version of eq. \eqref{eq:bte1} (i.e. with $q\equiv 0$ and b.c. \eqref{eq:bte3} or \eqref{eq:bte2} with
$\psi_{\text{in}} \equiv 0$) has a non-trivial non-negative solution. The free parameter is commonly placed in front of
the fission part and plays the role of an eigenvalue of eq. \eqref{eq:bte1} -- it is well known that the desired
solution is the unique eigenfunction associated with the simple eigenvalue of smallest modulus (see e.g. \cite[Chap.
XXI, \S3]{DautrayLions} or \cite{Sanchez3}). We will not, however, cover the critical eigenvalue problem here in any
more detail.

\subsection{Solvability}\label{sec:solvability}
Let $$
  \Hp(X) = \left\{\psi\in\Lp(X, \d{x}), \bomega\cdot\nabla\psi\in\Lp(X, \d{x}), \psi\vert_{\pX[-]}\in\Lp(\pX[-],
  \db)\right\}
$$ and
\begin{equation*}
  \begin{gathered}
    A\psi(\br,\bomega,E) = \bomega\cdot\nabla\psi(\br,\bomega,E),\\%[.85em]
    \Sigma_t\psi(\br,\bomega,E) = \sigma_t(\br,E)\psi(\br,\bomega,E),\\%[.75em]
    K\psi(\br,\bomega,E) = \intE[']{\Emin}{\Emax}{
            \intA[']{\kappa(\br,\bomega\cdot\bomega',E\sla E')\psi(\br,\bomega',E')}
          }.
  \end{gathered}
\end{equation*}
We shall call $A$, $\Sigma_t$, $K$ and $T = A + \Sigma_t - K$ the \textit{advection}, \textit{collision}, 
\textit{transfer} and \textit{transport} operator, respectively. Note that the operator $\Sigma_t - K$ is self-adjoint
 if and only if the kernel $\kappa$ of $K$ is symmetric in $E$ and $E'$; with the exception of the mono-energetic case, 
 this is generally not true (as already mentioned in the introduction). We also note that for piecewise smooth $\pV$ 
 and $1\leq p < \infty$, the traces $\psi\vert_{\pX[\pm]}\in\Lp(\pX[\pm], \db)$ are well defined for functions 
$$
  \psi\in \{\varphi\mid \varphi\in\Lp(X, \d{x}), \bomega\cdot\nabla\varphi\in\Lp(X, \d{x})\}
$$
and admit a continuous lifting (\cite[Chap. XXI]{DautrayLions}, \cite{Boulanouar1}). This lifting allows us to convert 
the problem with non-homogeneous boundary conditions \eqref{eq:bte2} to that with homogeneous ones, so we will deal only 
with the latter in this section. Existence results for the reflective or more general boundary conditions require 
special trace theorems, see \cite[Chap. XXI, Appendix of \S2]{DautrayLions}.


We will now follow the general idea of \cite[Chap. XXI, \S 2, Proposition 5]{DautrayLions} to prove the following

\begin{theorem}\label{thm1}
Assume that
\begin{enumerate}[label=(\alph*)]
	\item $\sigma_t \in \Lp[\infty](X)$, $\sigma_t \geq \mst > 0$ a.e. in $V\times[\Emin,\Emax]$,
	\item $\kappa \geq 0$ a.e. in $V\times\Sphere^2\times[\Emin,\Emax]^2$,
	\item $\displaystyle\esssup_{X} c(\br,\bomega,E) < 1$ where
	  \begin{equation}\label{eq:c}
	    c(\br,\bomega,E) := \frac{1}{\sigma_t(\br,E)}\int_{\Emin}^{\Emax}\int_{\Sphere} \kappa(\br,\bomega\cdot\bomega',
	    E\shortrightarrow E')\,\d{E'}\d{\bomega'}.
	  \end{equation}
\end{enumerate}
Then the fixed source, steady state neutron transport problem with vacuum boundary conditions 
\begin{equation*}
  \left\{
  \begin{aligned}
     &T\psi(\br,\bomega,E) = q(\br,\bomega,E),\\
     &\Dom{T} = \{\psi\in \Hp[1](X),\ \psi\vert_{\pX[-]} = 0\},
  \end{aligned}
  \right.
\end{equation*}
has a unique solution $\psi(\br,\bomega,E)\in\Dom{T}\subset\Lp[1](X)$ for any $q\in \Lp[1](X)$.
\end{theorem}
\begin{proof}
  The assumptions are equivalent to 
  \begin{enumerate}
    \item[(c')] $\displaystyle\qquad\qquad 0 \leq \intE[']{\Emin}{\Emax}{\intA[']{\kappa(\br,\bomega\cdot\bomega',E\shortrightarrow E')}} \leq \sigma_t(\br,E) - \mst$
  \end{enumerate}
  for a. e. $(\br,\bomega,E) \in X$ which implies that\footnote{the sharp brackets denote the duality pairing between 
  $\Lp[1](X)$ and $\Lp[\infty](X)$ via $\psi^\ast = \sgn \psi$}
  $$
    \begin{aligned}
    \langle \psi^\ast, K\psi \rangle &= \intX{\sgn \psi(\br,\bomega,E)\intE[']{\Emin}{\Emax}{\intA[']{\kappa(\br,\bomega\cdot\bomega',E\shortrightarrow E')\psi(\br,\bomega',E')}}}\\
    & \leq 
    \intE[']{\Emin}{\Emax}{\intA[']{\psi(\br,\bomega',E') \bigl[ \sigma_t(\br,E') - \mst \bigr]}}\\ &= \intX{\psi(\br,\bomega,E) \bigl[ \sigma_t(\br,E) - \mst \bigr]}.
    \end{aligned}
  $$
  Next,  
  $$
  \begin{aligned}
    \langle \psi^\ast, (\Sigma_t - \mst\Id)\psi\rangle &= \intX{\left[\sigma_t(\br,E) - \mst\right]\psi(\br,\bomega,E)\sgn \psi(\br,\bomega,E)} \\
    & = \intX{\left[\sigma_t(\br,E) - \mst\right]\abs{\psi(\br,\bomega,E)}} \geq \intX{\left[\sigma_t(\br,E) - \mst\right]\psi(\br,\bomega,E)}
  \end{aligned}
  $$
  and hence in summary
  $$
    \langle \psi^\ast, (\Sigma_t - K - \mst\Id)\psi\rangle \geq 0.
  $$
  Similarly
  $$
    \langle \psi^\ast, A\psi\rangle = \intX{(\bomega\cdot\nabla\psi)\sgn\psi} = \intX{\bomega\cdot\nabla\abs{\psi}}
    = \int_{\pX[+]}(\bomega\cdot\bn)\abs{\psi}\,\db \geq 0.
  $$
  using the Green's formula and vanishing inflow boundary conditions to obtain the last equality.
  Both operators are therefore m-accretive on $\Lp[1](X)$ (see \cite[Chap. XVII A]{DautrayLions}) and since the operator 
  $\Sigma_t - K - \mst\Id$ is bounded, also the whole operator 
  $$
    A + \Sigma_t - K - \mst\Id = T - \mst\Id
  $$ 
  is m-accretive on $\Lp[1](X)$. It then follows that 
  for every $\lambda > 0$, $\Id + \lambda (T - \mst\Id)$ is bijective from $D(T)$ onto $\Lp[1](X)$ and the conclusion 
  of the theorem is obtained for $\lambda = 1/\mst$.
\end{proof}
The value $c$ in Thm. \ref{thm1} has the physical meaning of the average number of neutrons emitted (with any velocity) 
after a neutron with velocity $(\bomega, E)$ collides with a nucleus at point $\br\in V$. Condition (c) thus expresses 
the requirement that the system be \textit{subcritical} in order for a steady solution in presence of external sources 
to be achieved\footnote{Note that $c$ slightly above $1$ may be allowed in practice due to the leakage of neutrons out 
of the system boundaries, which is not included in def. \eqref{eq:c}.}. Dautray and Lions in 
\cite[\S2, Chap. XXI]{DautrayLions} also outline the proof based on the same ideas in general $\Lp(X)$ spaces for 
$1 < p < \infty$, where in particular the case $p = 2$ is important\footnote{even though it does not lead to a direct 
physical interpretation as the $\Lp[1](X)$ setting} as the Hilbert space structure of $\Lp[2](X)$ allows to use richer 
set of mathematical tools to formulate practical solution methods (like those based on Fourier transform or the finite 
element methods). However, we need to add to assumptions (a-c) another assumption
\begin{enumerate}
  \item[(d)] ~~~~~~~~~~~~$\displaystyle 0 \leq \intE[']{\Emin}{\Emax}{\intA[']{\kappa(\br,\bomega\cdot\bomega',E\sla E')}} \leq \sigma_t(\br,E) - \mst$
\end{enumerate}
which does not have as plausible physical explanation as the assumption (c) (it could be interpreted as that beam of 
neutrons with any given velocity $(\bomega, E)$ is always attenuated by the collision with nuclei at any given point 
$\br\in V$ \footnote{i.e. the absorption rate dominates the rate at which new neutrons are introduced into the beam 
after collisions induced by all possible neutrons impinging on that point}). In Proposition 6, Dautray and Lions also 
state the existence of unique solution in $\Lp[\infty](X)$ for $q\in \Lp[\infty](X)$; the proof in this case is based on
 the inversion of the transport operator along characteristics\footnote{The integral form of the BTE, mentioned already 
 in the introduction, is obtained in this way.} and uses condition (d) instead of (c). Note that the approach based on 
 the characteristic form of the BTE has also been used by Sanchez in \cite{Sanchez3}, who proved the existence result in
  the cross-section weighted space $\Lp[1]_{\sigma_t}(X)$, using general boundary conditions. This seems to be an 
  alternative physically natural functional setting due to the definition of the reaction rate, eq. \eqref{eq:rr}; 
  moreover, assumption (a) may be relaxed by allowing $\sigma_t = 0$ in arbitrarily large regions 
  (the \textit{void regions}).

\subsection{Multigroup approximation}\label{sec:MG}
The above results hold also for the multigroup form of the transport equation (\cite[Chap. XXI, \S1]{DautrayLions}),
which is obtained by dividing the interval of neutron energies as follows\footnote{Note that the energy intervals
(groups) are numbered in a descending order, i.e. a group with larger index contains lower energies than a group with
lesser index.}:
$$
  [\Emin,\Emax] = [E^{G-{1/2}},E^{G+{1/2}}]\cup \ldots \cup [E^{g-1/2}, E^{g+1/2}] \cup \ldots \cup [E^{1/2},E^{3/2}],
$$ integrating \eqrefs{eq:bte1}{eq:bte3} over each energy group $[E^{g-1/2}, E^{g+1/2}]$ and defining appropriate
group-integrated quantities corresponding to the average group energy\footnote{A counting measure on the set
$\{E^G,\ldots,E^1\}$ is then used in the proofs instead of the continuous Lebesgue measure $\d{E}$ in the energy
integrals.}\footnote{We will always assume that the group parameters appearing in the equations (like macroscopic
cross-sections) are given and only remark that their actual calculation is an essential and highly non-trivial task.}
$E^g$. This conventional procedure leads to the following set of $G$ coupled neutron transport equations
\begin{equation}\label{eq:mg}
  \begin{gathered}
	\left\{
	  \begin{aligned}
      &T^\text{G}\{\psi^g(\br,\bomega)\} = \{q^g(\br,\bomega)\},\\
      &\Dom{T^\text{G}} = \bigl\{\{\psi^g\}\in \bigl[\Hp[1](X)\bigr]^G,\ \psi^g\vert_{\pX[-]} = 0,\ g = 1,\ldots,G\bigr\}
    \end{aligned}
  \right.\\[.2em]
    T^\text{G}\{\psi^g(\br,\bomega)\} = \left\{\left(A + \Sigma_r^g\right)\psi^g(\br,\bomega) - \summ{g'=1,g'\neq g}{G} K^{gg'}\psi^{g'}(\br,\bomega)\right\}\\
    \Sigma_r^g = \sigma_t^g(\br) - \intA[']{\kappa^{g\sla g}(\br,\bomega\cdot\bomega')},\quad K^{gg'} = \intA[']{\kappa^{g\sla g'}(\br,\bomega\cdot\bomega')}.
  \end{gathered}
\end{equation}
Note that it is customary to move the \textit{self-scattering} (diagonal) part of the transfer operator to the collision 
operator. Since the reactions in which energetic distribution of both the incoming and outgoing neutrons lies within the 
same group are included in both $\sigma_t$ and $\kappa$ (see Sec. \ref{sec:definitions}), this transformation makes 
$\Sigma_r^g\psi^g$ represent the actual rate of neutron removal from the group, while $K^{gg'}\psi^{g'}$ the rate of 
neutron addition to that group.

We conclude this chapter by recalling the traditional way of iterative solution of the multigroup system, known as 
\textit{source iteration}. Mathematically, it is a Gauss-Seidel iteration for the operator equation \eqref{eq:mg}:
$$
  \left(A + \Sigma_r^g - \summ{g'\leq g-1}{} K^{gg'}\right)\psi^g_{i+1} = \summ{g' \geq g+1}{} K^{gg'}\psi^{g'}_{i} + q^g,\quad g = 1,\ldots,G,\ \ i = 0,1,\ldots
$$
Now a mono-energetic transport problem in group $g$ has to be solved in each iteration, and only the advection operator 
$A$ spoils the self-adjointness of the formulation. 
%The following chapter is devoted to a particular self-adjoint approximation of the mono-energetic problem.



%  \chapter{The $\MCPN$ approximation}\label{sec:spn}

In this chapter, we start constructing an approximation to the monoenergetic neutron transport equation $T\psi - q = 0$
of the form
\begin{equation}\label{eq:MC-bte}
    \bomega\cdot\nabla\psi(\br,\bomega) + \sigma_t(\br)\psi(\br,\bomega) - \intA[']{\left[\sigma_s(\br,\bomega\cdot\bomega') + 
    \frac{\nu(\br)\sigma_f(\br)}{4\pi}\right]\psi(\br,\bomega')} - q(\br,\bomega) = 0.
\end{equation}
We will focus on the angular approximation here and assume that $\psi(\cdot,\bomega)\in\Lp[2](\Sphere)$
% in the discussion below (even though the final derivation in Sec. \ref{sec:mcpn} actually requires only $\Lp[1](\Sphere)$ integrability)
. 
Let us start by recalling the well established spherical harmonics method.

\section{$\PN$ approximation}\label{sec:SPH}
Being a square integrable function on $\Lp[2](\Sphere)$, $\psi(\cdot,\bomega)$ may be represented by a generalized
Fourier series expansion in terms of a complete orthonormal basis of this space. The $\PN$ method uses the basis of
tesseral spherical harmonic functions $\{\Y{n}{m}\}$ of degree $n\geq 0$ and order $-n \leq m \leq n$ (shortly
\textit{tesseral harmonics}). A finite method is obtained by the classical Galerkin approximation of the BTE in a
subspace $\Lp[2]_N(\Sphere)\subset\Lp[2](\Sphere)$ spanned by $\{\Y{n}{m}\}$,  $n \leq N$ \footnote{For subtle reasons
that could be found in the literature (\cite[Sec. 10.3.2]{Davison}, \cite[Sec. 9.6]{Stacey1}), only odd-$N$ $\PN$
approximations are usually used.} -- i.e. it consists of the expansion
\begin{equation}\label{eq:exp}
  \psi(\br,\bomega) \approx \suma[n]{0}{N}\suma[m]{-n}{n}\angmom{n}{m}(\br)\Y{n}{m}(\bomega)
\end{equation}
and orthogonal projection of the residual $T\psi - q$ onto $\Lp[2]_N(\Sphere)$ (similarly, the exact boundary conditions 
are projected onto a subspace of $\Lp[2](\pX[-])$ spanned by the complete set of odd-degree tesseral harmonics; 
the resulting conditions are known as the \textit{Marshak conditions}\footnote{We only make a remark that there is 
another form of approximate boundary conditions known as the Mark conditions. The relative merit of one over the other 
is not completely resolved so both are widely used. We choose the former as they are consistent with the Galerkin 
derivation of the $\PN$ equations.}).

The choice of the tesseral harmonics basis facilitates the realization of the above procedure. Besides the
orthonormality property, they obey a three-term recurrence rule that is used advantageously in expanding the advection
part
% (however, a closure relation such as that partial derivatives of $\angmom{N+1}{m}(\br)$ vanish for all relevant $m$
% becomes necessary)
and they also naturally appear in the \textit{addition theorem} for Legendre polynomials $P_n$:
\begin{equation}\label{eq:addition}
  P_n(\bomega\cdot\bomega') = \frac{4\pi}{2n+1}\suma[m]{-n}{n}\Y{n}{m}(\bomega)\Y{n}{m}(\bomega').
\end{equation}
This equality is used for simplifying the scattering part of the transfer operator, after its expansion in terms of 
Legendre polynomials (which form a complete orthonormal set in $\Lp[2]([-1,1])$):
\begin{equation}\label{eq:exp2}
  \sigma_s(\br,\bomega\cdot\bomega') \approx \suma[k]{0}{K}
      \frac{2k+1}{4\pi}\sigma_{sk}(\br)P_k(\bomega'\cdot\bomega)
\end{equation}
(where $K \leq N$ is usually called \textit{degree of scattering anisotropy}).
%where
%$\Sigma_{sk}(\br) = 2\pi \muint[_0]{\P{k}(\mu_0)\Sigma_s(\br, \mu_0)}$
%are the coefficients of the expansion of the scattering cross-section 
% ROTATIONAL INVARIANCE OF SUBSP.
Also, we directly obtain the important physical quantities as $\phi = \angmom{0}{0}$, $\mathbf{J} = \bigl[\angmom{1}{0},\angmom{1}{1},\angmom{1}{-1}\bigr]^T$.

\section{Other spherical harmonic approximations}
This standard procedure is described in various books on reactor physics and used in several computer codes
(\cite{MARC,Capilla,vanCriekingen1}). It results however in a complicated set of strongly coupled partial differential
equations without any clear structure that could be used for its simplification. This motivates the search for an
alternative set of expansion functions for \eqref{eq:exp} that would share with $\Y{n}{m}$ the important properties like
those mentioned above but have a more regular structure.

\subsection{Surface and solid spherical harmonics}
One possible such set arises by studying the linear combinations of tesseral harmonics of fixed degree. We will simplify
the notation by setting $\mu = \cost$, replace $\bomega$ with $(\mu, \azimuthal)$ where necessary and also use the
following relations between the velocity and direction vectors in Cartesian coordinates:
$$
  \bv = [v_x, v_y, v_z]^T = v\bomega,\quad \bomega = [\Omega_x, \Omega_y, \Omega_z]^T = \frac{\bv}{v}.
$$

\begin{definition}\label{defn:SSH}
  A general linear combination of the $2n + 1$ tesseral harmonics of degree $n$ is called a \textit{surface spherical 
  harmonic of degree $n$} and can be written as
  \begin{equation}\label{eq:exp3}
    Y_n(\mu,\azimuthal) = A_0 P_n(\mu) + \suma[m]{1}{n}\left[ A_m \cos(m\azimuthal)\P{n}{m}(\mu) + B_m \sin(m\azimuthal)\P{n}{m}(\mu)\right]
  \end{equation}
  where $\P{n}{m}$ are the \textit{associated Legendre polynomials} \cite[Chap. VI]{Byerly}. 
\end{definition}
Note that there is a one-to-one relationship between the coefficients in eq. \eqref{eq:exp} for fixed $n$ and the 
coefficients in \eqref{eq:exp3}.

Multiplying by $v^n$, we obtain the (regular) \textit{solid spherical harmonic}\footnote{Regular solid spherical
harmonics are one class of solutions of the Laplace equation $\lap Y = 0$ in spherical coordinates
$(v,\azimuthal,\polar)$ which vanish as $v\to 0$. The other are the irregular solid spherical harmonics, which have
singularity of the form $v^{-n-1}$ at the origin (\cite[Chap. VI]{Byerly}).}. Paper \cite{Ackroyd1} appears to be the
first where solid spherical harmonics were practically used to solve the neutron transport equation. Under the
assumption of isotropic scattering ($\sigma_{sk} = 0$ for $k \geq 1$ in \eqref{eq:exp2}) he derived for a homogeneous
region a set of coupled diffusion equations (called SH\PN) without any other requirement (unlike the classical
derivation of the $\SPN$ equations, which required certain assumptions about dimensionality or material properties, cf.
\cite{Larsen1}, \cite{Pomraning1}), together with heuristic boundary and interface conditions. As was shown in the
paper, the SH\PN~ equations reduced by simple substitutions to the set of $\SPN$ equations originally formulated by
Gelbard and it is interesting to note that this actually showed that the latter are within an isotropically scattering
homogeneous medium equivalent to the full solid harmonics expansion -- the same result for surface spherical harmonics
has been independently proven in \cite{Coppa1} and revisited recently (\cite{Coppa2,McClarren1}). Unfortunately, the
general treatment in \cite{Ackroyd1} is very technically involved and in author's opinion quite difficult to follow --
this may be the reason why the idea has not been picked up and possible research directions outlined in the paper's
conclusion not pursued.

\subsection{Cartesian tensors}\label{sec:tensors}
In order to develop a conceptually simpler and arguably also more useful approach, we need to recall some basic facts
about Cartesian tensors. We will use the convention that Greek subscripts represent axes of a Cartesian coordinate
system with unit vectors $\be_x$, $\be_y$, $\be_z$. We will also use the usual Einstein's summation convention which
implies summation over any index which appears twice in an indexed expression, for instance $$
  \AAA{\alpha\beta\gamma}\BBB{\beta\gamma} =
  \suma[\beta]{1}{3}\suma[\gamma]{1}{3}\AAA{\alpha\beta\gamma}\BBB{\beta\gamma} = \CCC{\alpha}.
$$

\begin{definition}
  An $n$-dimensional array $\AA$ of $3^n$ components $\AAA{\alpha_1\ldots\alpha_n}$ is called \textit{Cartesian tensor 
  of rank $n$} if it transforms as
  \begin{equation}\label{eq:tenstran}
    A^{(n)'}_{\alpha_1\ldots\alpha_n} = g_{\alpha_1\beta_1}\cdots g_{\alpha_n\beta_n}\AAA{\beta_1\ldots\beta_n}
  \end{equation}
  under the change of coordinate system $Oxyz \to Ox'y'z'$ by the action of an orthogonal matrix $\mathbf{G} = [g_{\alpha\beta}]$:
  $$
    \be_\alpha' = g_{\alpha\beta}\be_\beta.
  $$
\end{definition}
As we will only use the Cartesian tensors, we will henceforth omit the word Cartesian. Word ``tensor'' will also be used
 for a general \textit{tensor field}, components of which are functions -- like $\AA(\br)$. The transformation has then 
 the following form:
$$
  A^{(n)'}_{\alpha_1\ldots\alpha_n}(\br) = \AAA{\alpha_1\ldots\alpha_n}(\mathbf{G}^{-1}\br) = g_{\alpha_1\beta_1}\cdots g_{\alpha_n\beta_n}\AAA{\beta_1\ldots\beta_n}(\br).
$$
We will denote by $\eye$ the identity rank-2 tensor (matrix) and by $ \mathbb{O}^{(n)}$ the zero rank-n tensor.

Addition and subtraction of two tensors of same rank and multiplication of a tensor by a scalar are done component-wise. 
Multiplication of two tensors is defined as follows:
\begin{definition}
  Components of tensor $\CC[n+m] = \AA\otimes\BB[m]$ are given by
  $$
    \CCC[n+m]{\alpha_1\ldots\alpha_n\beta_1\ldots\beta_m} = \AAA{\alpha_1\ldots\alpha_n}\BBB{\beta_1\ldots\beta_m}.
  $$
\end{definition}

\begin{definition}
  $m$-th power of tensor $\AA$ is a tensor of rank $nm$ defined as
  $$
    \CC[nm] = \AA\otimes\AA\otimes\cdots\otimes\AA\quad \mbox{($m$-times)}
  $$
\end{definition}
We will mainly use powers of vectors and also consider the gradient operator as a vector
$$
  \nabla = \left[\pd{}{x}, \pd{}{y}, \pd{}{z}\right]^T
$$
(hence $\Delta = \lap = \nabla\otimes\nabla$). Also, to keep usual notation, $\nabla\AA \equiv \nabla\otimes\AA$.

\begin{definition}
  For $1\leq m \leq n$, the $m$-fold contraction of tensors $\AA$ and $\BB$ is a rank-($2n-2m$) tensor 
  $\CC[2n-2m] = \AA\underset{m}{\cdot}\BB$ with components
  $$
    \CCC[2n-2m]{\alpha_{1}\ldots\alpha_{n-m}\beta_1\ldots\beta_{n-m}} = 
    \AAA{\alpha_1\ldots\alpha_{n-m}\gamma_{n-m+1}\ldots\gamma_n}\BBB{\gamma_n\ldots\gamma_{n-m+1}\beta_{n-m}\ldots\beta_1}.
  $$
\end{definition}
Specially for $m = 1$, we get the standard inner product. 
In the case when $n = m$, we obtain the scalar $\AAA{\gamma_1\ldots\gamma_n}\BBB{\gamma_n\ldots\gamma_1}$ and 
suppress the index under the $\cdot$ sign to simplify the writing (it should be clear from the two operands and their 
rank that a total contraction over all their indices is intended). 

\begin{definition}
  For $m \leq \lfloor n/2 \rfloor$ (the integer part of $n/2$), the $m$-fold contraction of tensor $\AA$ (contraction 
  in $m$ index pairs) is a rank-($n-2m$) tensor with components
  $$
    \BBB[n-2m]{\alpha_{2m+1}\ldots\alpha_n} = \AAA{\alpha_1\alpha_1\ldots\alpha_m\alpha_m\alpha_{2m+1}\ldots\alpha_n}.
  $$
\end{definition}
When contracting in only 1 index pair, the $1$-fold contraction is simply called a contraction. The result has a special 
name:
\begin{definition}
  Contraction of tensor $\AA$ in one index pair is called \textit{trace} of $\AA$ in that pair:
  $$
    \trace \AA = \AAA[n-2]{\alpha_1\alpha_1\alpha_3\ldots\alpha_n}.
  $$ 
\end{definition}
If the trace of $\AA$ in one index pair vanishes, the tensor is said to be traceless in that pair. Obviously then, if 
the tensor does not depend on the order of indices, it is traceless in all index pairs.
\begin{definition}
  Tensor $\AA$ is \textit{totally symmetric} if it is invariant under any permutation $\perm{\alpha_1\ldots\alpha_n}$ 
  of its indices:
  $$
    \AAA{\alpha_1\ldots\alpha_n} = \AAA{\perm{\alpha_1\ldots\alpha_n}}.
  $$
\end{definition}

\begin{definition}
  Totally symmetric tensor whose trace in any index pair vanishes is called \textit{totally symmetric traceless tensor} 
  and its trace vanishes in all index pairs. We will denote totally symmetric traceless tensors as TST tensors.
\end{definition}

We can make any tensor of rank $n$ totally symmetric by applying the symmetrization operator $\Sym{\cdot}$:
\begin{definition}
  $\widetilde{\mathbb{A}}^{(n)} = \Sym\AA$ is a totally symmetric tensor with components
  $$
    \widetilde A^{(n)}_{\alpha_1\ldots\alpha_n} = \frac{1}{n!}\ \; \sum_{\mathclap{\perm{\alpha_1\ldots\alpha_n}}}\ \;\AAA{\alpha_1\ldots\alpha_n}.
  $$
\end{definition}  
As we will see shortly, there also exists an operator which makes a totally symmetric tensor totally symmetric and traceless.

\subsection{Maxwell-Cartesian spherical harmonics}\label{sec:MCSH}
It is well known that any solid spherical harmonic of degree $n$, when expressed in Cartesian coordinates: $S_n(\bv) =
S_n(v_x, v_y, v_z) = v^n Y_n(\bv/v)$, is a harmonic polynomial, homogeneous of degree $n$ (\cite[Art. 110]{Byerly}),
i.e.
$$
  \lapv S_n(\bv) = 0,\quad S_n(\lambda\bv) = \lambda^n S_n(\bv)
$$ (in fact, there is an alternative definition of solid spherical harmonic as a function satisfying these two
properties and of surface spherical harmonic as a restriction of such function to the unit sphere). The space of such
polynomials has dimension $2n + 1$\footnote{Homogeneous polynomial of degree $n$ has $\frac{(n+1)(n+2)}{2}$
coefficients; since $\lapv S_n(\bv)$ is a homogeneous polynomial of degree $n-2$ with $\frac{n(n-1)}{2}$ coefficients
expressed in terms of those of $S_n$, the condition $\lapv S_n(\bv) = 0$ reduces the number of independent coefficients
of $S_n$ to $\frac{(n+1)(n+2)}{2} - \frac{n(n-1)}{2} = 2n + 1$.}\label{pg:indepc} and is isomorphic to the space of TST
tensors of rank $n$ (\cite{Sam}).
% SEE AXLER HFT - L2(S_2) = direct sum of H_n(S_2)
In particular, when the components of such a tensor do not depend on $v_x$, $v_y$, $v_z$, (or $\Omega_x$, $\Omega_y$,
$\Omega_z$) they comprise coefficients of the solid (or surface) spherical harmonic and so also (in the
$\Sphere$-restricted case) the coefficients $\angmom{n}{m}(\br,E)$ in the expansion \eqref{eq:exp} (when transformed
back to spherical coordinates). Hence, by solving the set of tensorial equations with these tensors of ranks $n \leq N$
as unknowns, we could accomplish essentially the same as by solving the ordinary $\PN$ equations, but we could
additionally use the TST tensor character of the equations.

To construct this set, we will build on the ideas presented by Johnston (\cite{Johnston1}) and Applequist
(\cite{Applequist1}). Johnston arrived at a ``far more symmetric and compact'' (\cite[p. 1455]{Johnston1}) form of the
angularly discretized general Boltzmann-Vlasov equation by expanding the flow function in terms of
\textit{Maxwell-Cartesian spherical harmonic tensors} rather than the usual surface spherical harmonics. The term
``Maxwell-Cartesian spherical harmonic tensor'' has been coined by Applequist in \cite{Applequist2}, who in his older
paper \cite{Applequist1} showed a systematic way of obtaining a special TST tensor of any rank $n$ whose components are
spherical harmonics of degree $n$ in Cartesian frame of reference as defined by Maxwell in \cite[p. 160]{Maxwell}.
Specifically, in present notation, he showed that Maxwell's spherical harmonics based on Cartesian axes can be obtained
(up to a normalization constant) as components of $$
  \detr \bv^n \quad \mbox{(or $\detr \bomega^n$)}
$$ where $\detr$ is the so-called \textit{detracer operator} which projects a general totally symmetric tensor of rank
$n$ into the space of totally symmetric and traceless tensors of rank $n$ (see the note on p. 4311 of
\cite{Applequist1}). The result of applying the detracer operator on a totally symmetric tensor $\AA$ is a TST tensor
with components:
\begin{equation}\label{eq:detr}
  \detr \AAA{\alpha_1\ldots\alpha_n} = \suma[m]{0}{\lfloor n/2 \rfloor} \frac{(-1)^m(2n-2m-1)!! }{(2n-1)!!}\quad\sum_{\mathclap{\{\perm{\alpha_1\ldots\alpha_n}\}_{\text{d}}}}\quad\kron{\alpha_1}{\alpha_2}\cdots\kron{\alpha_{2m-1}}{\alpha_{2m}}\AAA{\beta_1\beta_1\ldots\beta_m\beta_m\alpha_{2m+1}\ldots\alpha_n},
\end{equation}
where $\{\perm{\alpha_1\ldots\alpha_n}\}_{\text{d}}$ denotes the set of all permutations that give distinct terms in 
the last sum (considering their total symmetry) and $(2n-1)!! = (2n-1)(2n-3)\cdots 3\cdot 1$ with $(-1)!! = 1$. Various 
properties and different explicit forms of the detracer operator are given in \cite[Sec. 5]{Applequist1}
%or, in a somewhat more transparent way, in \cite[Eqns. (36), (37)]{Vreoiju}
.

Being general surface spherical harmonics according to Def. \ref{defn:SSH}, Maxwell's surface spherical harmonics are
linear combination of tesseral harmonics and hence share many of the properties of the latter. An extensive presentation
of these properties is given in \cite{Applequist1,Applequist2} and we will recall some of them in further sections. We
will use the following form of Maxwell-Cartesian surface spherical harmonics of degree $n$ as components of rank-$n$ TST
tensors (called sometimes \textit{Maxwell-Cartesian tensors}):
\begin{equation}\label{eq:MCSH}
  \PP(\bomega) = \detr \bomega^n
\end{equation}
which differ from those used by Applequist in that they have coefficient $1$ at their first term but which coincide 
with those used by Johnston. 
 
We note that projection of $\PP$ along, say, $z$-axis (or any other because of the symmetry) yields (up to a
normalization factor) the Legendre polynomials:
\begin{equation}\label{eq:leg}
  \PP(\bomega)\contr\be_z^n = P^{(n)}_{\alpha_1\ldots\alpha_n}\kron{3}{\alpha_1}\cdots\kron{3}{\alpha_n} = P^{(n)}_{33\ldots3} = \frac{n!}{(2n-1)!!}P_n(\mu) \equiv C_n P_n(\mu)
\end{equation}
and hence components of $\PP$ could be explicitly obtained by extending the well-known formulas for Legendre 
polynomials, for instance
\begin{equation}\label{eq:conn1}
  \PP(\bomega) = \bomega^n - \frac{n(n-2)}{2(2n-1)}\Sym{\eye\otimes\bomega^{n-2}} + \frac{n(n-1)(n-2)(n-3)}{8(2n-1)(2n-3)}\Sym{\eye^2\otimes\bomega^{n-4}}\mp\ldots
\end{equation}
(\cite[Chap. VI]{Byerly}, \cite{Johnston1})
% or
%\begin{equation}\label{eq:conn2}
%  \PP[0](\bomega) = 1,\ \PP[n+1](\bomega) = \left[\bomega\otimes\PP[n](\bomega) - \tfrac{n^2}{4n^2-1}\eye\otimes\PP[n-1](\bomega)\right]_{\mbox{\scriptsize sym}},\ n = 1,2,\ldots.
%\end{equation}
First few Maxwell-Cartesian surface harmonics generated by this formula are shown in the table below together with the 
corresponding tesseral harmonics and Legendre polynomials.
\begin{center}
	\begin{tabular}{|A|B|C|D|}
	  \hline
		n & $\YY(\bomega)\propto$ & $\PP(\bomega)$ & $C_n P_n(\mu)$ \nl
		0 & 1 & 1 & 1\nl
		1 & $\Omega_x$, $\Omega_y, \Omega_z$ & $\Omega_x$, $\Omega_y$, $\Omega_z$ & $\mu$ \nl
		2 & $-\Omega_x^2-\Omega_y^2 + 2\Omega_z^2$, 
		    $\Omega_y\Omega_z$, 
		    $\Omega_z\Omega_x$,
		    $\Omega_x\Omega_y$,
		    $\Omega_x^2 - \Omega_y^2$ & $\Omega_x^2-\tfrac13$, 
		                                $\Omega_y^2-\tfrac13$, 
		                                $\Omega_z^2-\tfrac13$, 
		                                $\Omega_x\Omega_y$, 
		                                $\Omega_x\Omega_z$, 
		                                $\Omega_y\Omega_z$ & $\mu^2 - \tfrac13$ \nl
	\end{tabular}
	\captionof{table}{Tesseral and Maxwell-Cartesian surface spherical harmonics and Legendre polynomials up to degree 
	$n = 2$.}
	\label{tab:harmonics}
\end{center}
 The symmetric nature of the Maxwell-Cartesian surface harmonics as opposed to the tesseral harmonics can be seen from 
 the third row in the table. We also note that since there are $\frac{(n+1)(n+2)}{2}$ distinct components in a general 
 totally symmetric tensor of rank $n$, the number of Maxwell-Cartesian harmonics is greater than the number of tesseral 
 harmonics of the same degree, which is ($2n + 1$) for $n \geq 2$. However, due to the $\frac{n(n-1)}{2}$ conditions 
 arising from the tracelessness property, there are just $2n + 1$ independent components in $\PP$ (see also the footnote 
 on pg. \pageref{pg:indepc}). This also indicates the fact that unlike the tesseral harmonics with same degree but
 different orders, components of $\PP$ are not all $\Lp[2](\Sphere)$-orthogonal. Nevertheless, Maxwell-Cartesian surface
  tensors of different degrees are orthogonal in the following sense:
\begin{equation}\label{eq:OG}
  \intA{\PP[n](\bomega) \otimes \PP[m](\bomega)} = \mathbb{O}^{m+n},\quad n\neq m.
\end{equation}

The Maxwell-Cartesian tensors have been actively used for solving various electro-magnetics and quantum-mechanical
problems (see the references in \cite{Applequist2}), but there is apparently only one attempt to bring them to the field
of neutron transport. In a recent article \cite{Coppa3}, Coppa re-derived the Maxwell-Cartesian surface spherical
harmonic tensors from their analogy with Legendre polynomials (see above) and the requirement that they satisfy the
addition theorem of the form \eqref{eq:addition}. He then multiplied the BTE with these tensors and integrated over the
sphere to obtain a system of tensorial equations, each of which being itself a symmetric set of PDE's determining the
components of the flux moment tensor defined by
\begin{equation*}
  \psi^{(n)}(\br) = \intA{\psi(\br,\bomega)\PP[n](\bomega)}.\tag{\cite[Eqn. (42)]{Coppa3}}
\end{equation*}
Note that this definition implies that the flux moment tensors $\psi^{(n)}$ are totally symmetric and traceless. 
However, only the symmetry property appears to be recognized in \cite{Coppa3}. Using rather complicated tensor relations 
(which are not fully derived in the paper and only conjectured to hold for general $n$), the system is constructed again 
by analogy with the procedure leading to the $\PN$ system. This could be misleading, however. In particular, there is 
no indication of why the moments obtained by solving the set could be interpreted as coefficients in an expansion of 
angular flux. Note that in the usual $\PN$ method where the angular flux is expanded into a series of tesseral 
harmonics, the key ingredient that allows us to deduce 
%that allow us to deduce from the residual projection condition
%$$
%  \suma[n']{0}{\infty}\suma[m']{-n'}{n'}\Y{n'}{m'}(\bomega)\intA{\left[T\psi(\br,\bomega) - q(\br,\bomega)\right] \Y{n}{m}(\bomega)} = 0
%$$
%that 
%$$
%  \intA{\left[T\psi(\br,\bomega) - q(\br,\bomega)\right] \Y{n}{m}(\bomega)} = 0\quad \forall n,m
%$$
%(which leads to the moment equations analogous to \cite[Eqn. (48)]{Coppa3}) 
the system of equations analogous to \cite[Eqn. (48)]{Coppa3}) for the coefficients in the expansion is orthogonality 
(and hence linear independence) of all $\Y{n}{m}$. We do not have this property in the set of Maxwell-Cartesian 
harmonics (which seems to be neglected in \cite{Coppa3}). This issue however, can be elegantly resolved by taking into 
account the tracelessness property of the Maxwell-Cartesian tensors. We will demonstrate it by combining the results of 
Johnston (\cite{Johnston1}) and Applequist (\cite{Applequist1}). Surprisingly, this more fundamental and systematic 
approach leads to a set of equations equivalent with \cite[Eqn. (42)]{Coppa3} (explanation of why this is so deserves 
further investigation). We will also see a clear connection to the 1D $\PN$ equations (based on the expansion of angular 
flux in Legendre polynomials) which could be expected from the formal equivalence of the Maxwell-Cartesian tensors and 
Legendre polynomials mentioned above (\eqref{eq:conn1}, 
%\eqref{eq:conn2}, 
Tab. \ref{tab:harmonics})\footnote{In \cite{Coppa3}, only the equivalence with $\PN[1]$ and $\SPN[2]$ equations have 
been shown by rather laborious manipulations with the equations for $l = 0,1$ and $l = 0,1,2$, respectively.}.

\section{Derivation of the $\MCPN$ approximation}\label{sec:mcpn}
\subsection{First attempts}
Using the results of \cite[Art. 114]{Byerly}, the expansion of angular neutron flux in terms of surface spherical 
harmonics could be written as
\begin{equation}\label{eq:MC-exp1}
  \psi(\br,\bomega) = \suma[n]{0}{\infty}\frac{2n+1}{4\pi}\intA{\psi(\br,\bomega')P_n(\bomega\cdot\bomega')}.
\end{equation}
As shown in \cite[Sec. 7, Corollary II]{Applequist1}, Maxwell-Cartesian tensors satisfy the following form of 
addition theorem:
\begin{equation}\label{eq:MC-addition}
  \PP(\bomega)\contr\PP(\bomega') = C_n P_n(\bomega\cdot\bomega')
\end{equation}
($C_n$ defined in \eqref{eq:leg}). Combining the two results, we obtain
\begin{equation}\label{eq:MC-exp2}
  \psi(\br,\bomega) = \suma[n]{0}{\infty}\frac{2n+1}{4\pi C_n}\left[\intA{\psi(\br,\bomega')\PP(\bomega')}\right]\contr\PP(\bomega).
\end{equation}
If we now define the \textit{$n$-th angular flux moment tensor} by
\begin{equation}\label{eq:flux-moment}
  \fmom(\br) := \frac{2n+1}{4\pi C_n}\intA{\psi(\br,\bomega)\PP(\bomega)}
\end{equation}
equation \eqref{eq:MC-exp2} becomes
\begin{equation}\label{eq:MC-exp3}
  \psi(\br,\bomega) = \suma[n]{0}{\infty}\fmom(\br)\contr\PP(\bomega).
\end{equation}
Note that we have, like before, $\phi = \fmom[0]$, $\mathbf{J} = \bigl[\fmom[1]_x,\fmom[1]_y,\fmom[1]_z\bigr]^T$.

Although equation \eqref{eq:MC-exp3} with \eqref{eq:flux-moment} looks like a (generalized) Fourier series, we should 
keep in mind that
$$
  \intA{\PP[\ell](\bomega)\otimes\suma[n]{0}{\infty}\fmom(\br)\contr\PP(\bomega)} \neq \fmom[\ell]
$$
in general and regard \eqref{eq:flux-moment} as definition.

We analogously expand also the volumetric source term, leading to 
\begin{equation}\label{eq:MC-exp_src}
  q(\br,\bomega) = \suma[n]{0}{\infty}\qmom(\br)\contr\PP(\bomega)
\end{equation}
with
\begin{equation*}
  \qmom(\br) := \frac{2n+1}{4\pi C_n}\intA{q(\br,\bomega)\PP(\bomega)}.
\end{equation*}

To find the relations that must be satisfied by the angular expansion moments $\fmom(\br)$ in order for
\eqref{eq:MC-exp3} (or equivalently \eqref{eq:MC-exp1}) to be the solution of the BTE \eqref{eq:MC-bte}, we insert the
expansion \eqref{eq:MC-exp3} into \eqref{eq:MC-bte} (with source term represented according to \eqref{eq:MC-exp_src}).
Let us first look at the transfer part.
Applying eq. \eqref{eq:MC-addition} in the expansion \eqref{eq:exp2}, we immediately simplify the scattering
term\footnote{As in \eqref{eq:MC-exp1}, the expansion is taken exactly with $K\to\infty$ so that we can simply write
equalities between the $\Lp[2](\bomega)$ functions and their generalized Fourier series and use a single index $n =
0,1,\ldots$ in both expansions. In a practical calculation the expansion in \eqref{eq:MC-exp1} is truncated at finite
degree $N$ and the scattering cross-section moments $\sigma_{sk}$ are typically available only up to a certain
anisotropy degree $K \leq N$; in that case, however, we may always set $\sigma_{sn} = \sigma_{sk}$ for $n \leq K$ and
$\sigma_{sn} = 0$ otherwise.}:
\begin{equation*}%\label{eq:scat}
  \begin{aligned}
  \intA[']{\sigma_s(\br,\bomega\cdot\bomega')\psi(\br,\bomega')} &= 
    \suma[n]{0}{\infty}
    \frac{2n+1}{4\pi}\sigma_{sn}(\br)\intA[']{P_n(\bomega'\cdot\bomega)\psi(\br,\bomega')}\\
    & = \suma[n]{0}{\infty}
    \frac{2n+1}{4\pi C_n}\sigma_{sn}(\br)\left[\intA[']{\PP(\bomega')\psi(\br,\bomega')}\right]\contr\PP(\bomega)\\
    & = \suma[n]{0}{\infty}
    \sigma_{sn}(\br)\fmom(\br)\contr\PP(\bomega).
    \end{aligned}
\end{equation*}
Because of the orthogonality relation \eqref{eq:OG}, the fission part will have the following form:
\begin{equation}%\label{eq:scat}
  \intA[']{\frac{\nu(\br)\sigma_f(\br)}{4\pi}\psi(\br,\bomega')} = 
    \suma[n]{0}{\infty}
    \frac{\nu(\br)\sigma_f(\br)}{4\pi}\fmom(\br)\contr\intA[']{\PP(\bomega')} = \nu(\br)\sigma_f(\br)\phi(\br).
\end{equation}
Therefore, by inserting \eqref{eq:MC-exp3} into \eqref{eq:MC-bte} and using these results we obtain
\begin{equation}\label{eq:MC-bte2}
  \suma[n]{0}{\infty}\left[\bomega\cdot\nabla\fmom + \sigma_t\fmom 
  - \sigma_{sn}\fmom[n] - \kron{n}{0}\nu\sigma_f\phi - \qmom\right]\contr\PP(\bomega) = 0
\end{equation}
where each term in the brackets is dependent only on $\br$ (which is omitted for brevity).

Formulation of the moment equations for $\fmom(\br)$ is now hampered by linear dependence among certain functions in
each $\PP(\bomega)$ (so that we cannot deduce from \eqref{eq:MC-bte2} that for each $n$, all components of the
expression in brackets must vanish), as well as by the advection term which still contains $\bomega$. Both issues may be
however overcome by the so-called \textit{detracer exchange theorem} (\cite[Sec. 5.2]{Applequist2}).
\begin{theorem}\label{thm:detex}
  If $\AA$ and $\BB$ are totally symmetric tensors of rank $n$, then
  $$
    \AA\contr\detr \BB = \BB\contr\detr \AA.
  $$
\end{theorem}
\begin{proof}
  The theorem easily follows from the definitions of the detracer operator \eqref{eq:detr} and tensor contraction.
\end{proof}
Since $\fmom(\br)$ is by definition totally symmetric and traceless and $\bomega^{n}$ totally symmetric, using Theorem 
\ref{thm:detex} with $\AA \equiv \fmom(\br)$ and $\BB \equiv \bomega^{n}$ and the definition \eqref{eq:MCSH} of 
Maxwell-Cartesian tensors shows that\footnote{Recall that the detracer is a projection operator into the space of 
TST tensors and thus leaves the already TST tensor unchanged.} 
$$
  \fmom(\br)\contr\PP(\bomega) = \fmom(\br)\contr\bomega^n
$$
and the expansion \eqref{eq:MC-exp3} is thus equivalent to a power series in $\bomega$:
$$
  \psi(\br,\bomega) = \suma[n]{0}{\infty}\fmom(\br)\contr\bomega^n
$$
(similarly for \eqref{eq:MC-exp_src}). The advection term therefore simplifies as follows:
$$
  (\bomega\cdot\nabla)\fmom\contr\bomega^n = \left.
  \pw{\Omega_\beta\gradcomp\beta\fmom_{\alpha_0,\ldots,\alpha_{n-1}}\Omega_{\alpha_{n-1}}\cdots\Omega_{\alpha_0}}{if $n \geq 1$}
     {\Omega_\beta\gradcomp\beta\phi}{if $n = 0$}\right\}
   = \nabla\fmom\contr[n+1]\bomega^{n+1}
$$
for $n = 0,1,\ldots$ and by reindexing:
$$
  (\bomega\cdot\nabla)\fmom[n-1]\contr\bomega^{n-1} = \nabla\fmom[n-1]\contr[n]\bomega^{n}
$$
for $n = 1,2,\ldots$. We can therefore rewrite eq. \eqref{eq:MC-bte2} as
\begin{equation}\label{eq:MC-bte3}
  \suma[n]{0}{\infty}\left[\nabla\fmom[n-1] + \sigma_t\fmom
  - \sigma_{sn}\fmom - \kron{n}{0}\nu\sigma_f\phi - \qmom\right]\contr\bomega^n = 0
\end{equation}
with the term with negative rank discarded. Although eq. \eqref{eq:MC-bte3} closely resembles a vanishing linear 
combination of monomials
\footnote{Since we will not need to distinguish between the position vector $\br$ and the velocity vector $\bv$ here, 
we use the standard Cartesian coordinates $[x,y,z]$ for the latter as well.}
$$
  \bv^n = x^{n_1}y^{n_2}z^{n_3},\quad n_1 + n_2 + n_3 = n
$$
whose linear independence would imply that the expression in brackets must be a zero tensor for each $n$, by examining 
the not fully sensible equations that we would obtain in this way we might suspect that this is not the case. 
The reason is given in the following section.

\subsection{Sufficient conditions for deriving the moment equations}
When monomials $\bv^n$ of different degrees are restricted to the unit sphere, there appear nontrivial but vanishing
linear combinations of them (consider e.g. $\Omega_x^2 + \Omega_y^2 + \Omega_z^2 - 1$). However, when only the
components of $\bomega^n$ for a \textsl{single} $n$ are considered, they \textsl{are} linearly independent according to
the following lemma.
\begin{lemma}
  If $\AA$ is a totally symmetric tensor of rank $n$, independent of $\bomega$, and
  \begin{equation}\label{eq:lemma1}
    \AA\contr\bomega^n = 0,
  \end{equation}
  then $\AA = \mathbb{O}^{(n)}$.
\end{lemma}
\begin{proof}
  If $\AA$ is a totally symmetric tensor, then 
  $$
    \AA\contr\bomega^n = \frac{1}{v^n}\AA\contr\bv^n
  $$
  represents a linear combination of monomials of degree $n$,  with coefficients being the components of $\AA$. By 
  linear independence of these monomials we conclude that if the equality \eqref{eq:lemma1} holds, then all components 
  of $\AA$ must vanish.
\end{proof}

One possibility how to extend the above result to an arbitrary linear combination of direction tensors with different
ranks is to require the coefficient tensors to be not only totally symmetric, but also traceless.

\begin{theorem}\label{thm:MC-suffcond}
  If $\AA$ is a TST tensor independent of $\bomega$ for all $n = 0,1,\ldots$ and
  $$
    \suma[n]{0}{\infty}\AA\contr\bomega^n = 0,
  $$
  then $\AA = \mathbb{O}^{(n)}$ for all $n$.
\end{theorem}
\begin{proof}
  As discussed in Sec. \ref{sec:MCSH}, any TST tensor $\AA$ uniquely identifies a surface spherical harmonic of degree 
  $n$ by 
  $$
    Y_n(\bomega) = \AA\contr\bomega^n.
  $$
  Since surface spherical harmonics of different degrees are linearly independent,
  $$
    \suma[n]{0}{\infty} Y_n(\bomega) = 0\quad \Rightarrow \quad Y_n(\bomega) = 0 \ \ \forall{n}.
  $$
  The conclusion follows from the above lemma.
\end{proof}

\subsection{Final set of moment equations}
Using Thm. \ref{thm:MC-suffcond}, we see that we can obtain from \eqref{eq:MC-bte3} a set of moment equations if we
ensure that the expression in square brackets is a TST tensor for each $n$. Since $\fmom$ and $\qmom$ are TST by
definition, we only have to symmetrize and detrace the advection terms. We need to be careful, however, not to change
the original equation. This can be done by a clever rearranging of the terms in the sum.

First note that thanks to the symmetry of $\bomega^n$ we obviously have $$
  \nabla\fmom[n-1]\contr\bomega^n = \Sym{\nabla\fmom[n-1]}\contr\bomega^n.
$$ The symmetrized advection terms can be put into an equivalent TST form by $$
  \Sym{\nabla\fmom[n-1]} = \detr\Sym{\nabla\fmom[n-1]} - \left(\detr\Sym{\nabla\fmom[n-1]} -
  \Sym{\nabla\fmom[n-1]}\right).
$$ The result of the detracer operation on $\Sym{\nabla\fmom[n-1]}$ can be seen by using $$
  \AAA{\alpha_1\ldots\alpha_n} = \Sym{\gradcomp{\alpha_1}\fmom[n-1]_{\alpha_2\ldots\alpha_n}},\quad n \geq 2
$$ in eq. \eqref{eq:detr}. Taking into account the tracelessness of $\fmom[n-1]$, the only nonzero terms in the outer
sum will be those with $m = 0$ and $m = 1$. For $m = 0$, we obtain just the original tensor $\AA =
\Sym{\nabla\fmom[n-1]}$, while for $m = 1$, we obtain the inner sum over terms
\begin{equation*}
  \Sym{\kron{\alpha_1}{\alpha_2}\gradcomp{\beta}\fmom[n-1]_{\beta\alpha_3\ldots\alpha_n}},\quad
  \Sym{\kron{\alpha_1}{\alpha_3}\gradcomp{\beta}\fmom[n-1]_{\alpha_2\beta\alpha_4\ldots\alpha_n}},\quad\mbox{etc.},
\end{equation*}
that is $n-1$ times $\Sym{\eye\otimes\nabla\cdot\fmom[n-1]}$. We therefore get
\begin{equation}
\begin{aligned}\label{eq:MC-split0}
  \nabla\fmom[n-1]\contr\bomega^n &= \Sym{\nabla\fmom[n-1]}\contr\bomega^n\\
    & = \Sym{\nabla\fmom[n-1] - \frac{n-1}{2n-1}\eye\otimes\nabla\cdot\fmom[n-1] + \frac{n-1}{2n-1}\eye\otimes\nabla\cdot\fmom[n-1]}\contr\bomega^n\\
    & = \Sym{\nabla\fmom[n-1] - \frac{n-1}{2n-1}\eye\otimes\nabla\cdot\fmom[n-1]}\contr\bomega^n + \frac{n-1}{2n-1}\left(\nabla\cdot\fmom[n-1]\right)\contr\bomega^{n-2}
\end{aligned}
\end{equation}
where the last equality comes from the realization that $\eye\cdot\bomega^2 = \bomega\cdot\bomega = 1$ and 
$\nabla\cdot\fmom[n-1]$ is totally symmetric since $\fmom[n-1]$ is. 
Note that we have actually obtained that
\begin{equation}\label{eq:MC-split}
  \nabla\fmom[n-1]\contr\bomega^n = \detr\Sym{\nabla\fmom[n-1]}\contr\bomega^n + 
  \frac{n-1}{2n-1}\left(\nabla\cdot\fmom[n-1]\right)\contr\bomega^{n-2}.
\end{equation}
Eq. \eqref{eq:MC-split} shows that even though $\nabla\fmom[n-1]$ by itself is neither symmetric nor traceless, its 
contraction with $\bomega^n$ can be split into a sum of two contractions of TST tensors (total symmetry and 
tracelessness of the latter follows from the fact that $\fmom[n-1]$ is TST). By introducing this splitting into eq. 
\eqref{eq:MC-bte3} and regrouping the sum by $\bomega^n$, we finally obtain
\begin{equation}\label{eq:MC-bte4}
\begin{multlined}
  \suma[n]{0}{\infty}\left\{\Sym{\nabla\fmom[n-1] - \frac{n-1}{2n-1}\eye\otimes\nabla\cdot\fmom[n-1]} + \frac{n+1}{2n+3}\nabla\cdot\fmom[n+1]+\right.\\
  \left.\vphantom{\Sym{\nabla\fmom[n-1] - \frac{n-1}{2n-1}\eye\otimes\nabla\cdot\fmom[n-1]} + \frac{n+1}{2n+3}\nabla\cdot\fmom[n+1]}
     + \sigma_t\fmom - \sigma_{sn}\fmom - \kron{n}{0}\nu\sigma_f\phi - \qmom
  \right\}\contr\bomega^n = 0
\end{multlined}
\end{equation}
(again with the non-sensical tensors with negative ranks discarded).

Equation \eqref{eq:MC-bte4} is completely equivalent to eq. \eqref{eq:MC-bte3} or \eqref{eq:MC-bte2}, but regrouped to a
form of a vanishing linear combination of $\bomega^n$ with TST coefficient tensors for each $n$. Theorem
\ref{thm:MC-suffcond} therefore applies and we get the desired set of first-order partial differential equations for the
angular flux moments:
\begin{equation}\label{eq:MC-set}
  \frac{n+1}{2n+3}\nabla\cdot\fmom[n+1] + \Sym{\nabla\fmom[n-1] - \frac{n-1}{2n-1}\eye\otimes\nabla\cdot\fmom[n-1]} + 
  \sigma_t\fmom - \sigma_{sn}\fmom - \kron{n}{0}\nu\sigma_f\phi = \qmom
\end{equation}
for each $n = 0,1,\ldots$. 

\subsection{The $\MCPN$ approximation}
Note that in passing to a finite approximation by the simplest closure $\nabla\cdot\fmom[N+1] \equiv \mathbb{O}^{N}$ for
some $N \geq 0$, we do not spoil the TST character of the $N$-th coefficient tensor (in view of \eqref{eq:MC-split0},
this closure actually means that we are neglecting the nonzero traces of $\nabla\fmom[N-1]$).
The set \eqref{eq:MC-set} for $n \leq N$ may be regarded as an alternative to the ordinary $\PN$ equations and because
its solution represents the expansion of angular flux into Maxwell-Cartesian surface spherical harmonics of degrees up
to $N$, it will be called \textit{$\MCPN$ approximation}. We now explicitly state the first four equations in the
$\MCPN[3]$ set. We put $$\Sa{n} = \sigma_t - \sigma_{sn} - \kron{n}{0}\nu\sigma_f$$ and also include the $\otimes$
symbol whenever multiplication of tensors of rank $\geq 1$ occurs and finally obtain

\begin{equation}\label{eq:MCP3}
\begin{aligned}
    \tfrac{1}{3}\nabla\cdot\fmom[1] 
      +   {\Sigma}_{0}\phi
      &=  {q}^{(0)} \\[.2em]
    \tfrac{2}{5} \nabla\cdot\fmom[2] + \Sym{\nabla\otimes\phi} 
      +   {\Sigma}_{1}\fmom[1] 
      &=  {q}^{(1)}  \\[.2em]
    \tfrac{3}{7} \nabla\cdot\fmom[3] + \Sym{\nabla\otimes\fmom[1]} - \tfrac{1}{3} \Sym{\eye\otimes\nabla\cdot\fmom[1]}
      +   {\Sigma}_{2}\fmom[2]
      &=  {q}^{(2)}  \\[.2em]
                                       \Sym{\nabla\otimes\fmom[2]} - \tfrac{2}{5} \Sym{\eye\otimes\nabla\cdot\fmom[2]}
      +   {\Sigma}_{3}\fmom[3]
      &=  {q}^{(3)}
    \end{aligned}
\end{equation}

We see that the $\MCPN[1]$ set is equivalent to the $\PN[1]$ set. For higher $N$, we have more equations in the $\MCPN$
set than in the corresponding $\PN$ set, but because the coefficient tensors are bound to be TST, we can choose for each
$n$ any $2n+1$ equations and solve them for the corresponding $2n+1$ moments (which are the coefficients in
\eqref{eq:MC-exp3} at the functions forming an irreducible linearly independent subset of $\PP(\bomega)$) and express
the remaining ones as their linear combinations.

More important however appears the potential of using the symmetric and traceless structure of the $\MCPN$ equations to
provide new perspectives on some other widely used approximations or to create new ones. For instance, by projecting
each tensor along a chosen axis (say $z$), we obtain one dimensional equations which are exactly the same as the 1D
$\PN$ equations (since both symmetrized terms in every $n \geq 2$ $\MCPN$ equation become equal to
$\der{\fmom[n-1]_z}{z}$) if we take into account the definition of moments in the classical $\PN$ equations and multiply
each $\fmom[n]_z$ accordingly by a factor ${2n+1}$. This indicates the possibility to investigate the original ad-hoc
derivation of the $\SPN$ equation by formal extension of the 1D $\PN$ equations into 3D in the current tensorial
framework. Similarly, by multiplying each $\fmom[n]$ by $\frac{n!}{(2n-1)!!}$, we arrive at the system derived by Coppa
(\cite{Coppa3}).

There is also an interesting link to an old article of Selengut (\cite{Selengut}) in which the full multidimensional
$\PN[3]$ solution is obtained by solving a set of two coupled diffusion equations\footnote{much like in the $\SPN[3]$
approximation, but apparently without restrictions on dimensionality or cross-sections other than the usual isotropic
scattering and volumetric source assumptions} with special interface conditions in presence of multiple heterogeneous
regions. Selengut's derivation of the set is however quite puzzling (see also commentary in \cite[Sec. 5.2]{McClarren2})
and his equations have never been either analyzed or at least numerically tested. Assuming sufficient differentiability
of the angular flux moments $\fmom$ in eq. \eqref{eq:MCP3}, isotropic volumetric sources ($q^{(n)} = 0$ for $n\geq 1$)
and no void regions (i.e., all $\Sigma_{rn} > 0$), one can easily combine all equations into a single equation of the
form
\begin{equation}\label{eq:mom2}
  \fmom[2] = F(\phi, \fmom[2]) + \frac{1}{3\Sigma_2}\eye\Bigl(q^{(0)} - \Sigma_0\phi\Bigr)
\end{equation}
with $F$ being a differential operator. The right hand side of \eqref{eq:mom2}, however, is merely a symmetric tensqor
(matrix) with non-zero trace, and \eqref{eq:mom2} is thus not compatible with our definition of $\fmom[2]$. By 
adding the \textit{compatibility condition} $$
  \trace\left[F(\phi, \fmom[2]) + \frac{1}{3\Sigma_2}\eye\Bigl(q^{(0)} - \Sigma_0\phi\Bigr)\right] = 0,
$$
however, one can eliminate the components of $\fmom[2]$ and get the following equation for scalar flux
\footnote{This has been done so far only for the $\MCPN[3]$ set using the computer algebra system Mathematica, but a 
general procedure can be almost certainly formulated, leading to scalar flux equation with $\nabla^6$ term for
 $\MCPN[5]$ and so forth.}
\begin{equation}\label{eq:scfl}
  \frac{3}{35\Sigma_1\Sigma_2\Sigma_3}\nabla^4\phi + \left(\frac{4}{15\Sigma_1\Sigma_2} + 
  \frac{9}{35\Sigma_2\Sigma_3}\right)\nabla^2(q_0 - \Sigma_{0}\phi) - \frac{1}{3\Sigma_1}\nabla^2\phi
	  = q_0 - \Sigma_{0}\phi,
\end{equation}
which in the case of isotropic scattering reduces exactly to Selengut's Eqn. (5). Further investigation of this 
connection may be useful as it could shed some light on Selengut's derivation of interface conditions or lead to a
 new one (including the boundary conditions).





%: ----------------------- bibliography ------------------------

%\begin{multicols}{2} % \begin{multicols}{ # columns}[ header text][ space]
%\begin{tiny} % tiny(5) < scriptsize(7) < footnotesize(8) < small (9)

%\addcontentsline{toc}{chapter}{Bibliography}
%\bibliographystyle{alphaurl}
\bibliographystyle{plainnat} % Title is link if provided
\renewcommand{\bibname}{Literature} % changes the header; default: Bibliography

\bibliography{literature} % adjust this to fit your BibTex file

%\end{tiny}
%\end{multicols}

%: Declaration of originality
%\begin{declaration}

I hereby declare that this doctoral thesis is my own work, unless clearly stated otherwise.
\vspace{1cm}

\begin{flushright}
 \ldots\ldots\ldots\ldots\ldots\ldots\ldots\ldots\ldots\\
 Milan Hanu{\v s}~~~~~~~~~~
\end{flushright}
	
\end{declaration}


\end{document}
